%! Author = Profesor
%! Date = 06/06/2024


\section{Ejercicios}

\subsection{Ejercicios generales}

\begin{exercise}
%? \problema{\marcador{Ejercicio 33, página 191}}
    Se dispone de tres recipientes en los cuales tiene lugar una reacción química gaseosa que llega al equilibrio:
    \begin{equation*}
        a)\ \formula{A + B <=> C} \qquad b)\ \formula{D  <=> E + F} \qquad c)\ \formula{G + H <=> L + M}
    \end{equation*}
    Hay que elegir un único procedimiento que, en los tres casos, haga aumentar la proporción de los productos.
    ¿Sería alguno de los siguientes?: Aumento de la presión total.
    Disminución de la concentración de reactivos.
    Presencia de un catalizador.
    \begin{solu}
        \begin{itemize}
            \item El aumento de presión no, porque, en el segundo caso, desplazaría la reacción hacia la izquierda.
            \item La disminución de la concentración de reactivos en ninguna producirá un aumento de productos.
            \item La presencia de un catalizador puede aumentar la velocidad de una reacción pero en ningún caso variar las proporciones de productos y reactivos en el equilibrio.
        \end{itemize}
        Un procedimiento general sería aumentar las concentraciones de los reactivos y/o disminuir las de
        los productos.
    \end{solu}
\end{exercise}

\begin{exercise} %?\problema{\marcador{Ejercicio 34, página 191}}
    Para que una reacción alcance el equilibrio, ¿es necesario que se realice en un recipiente cerrado?
    \begin{solu}
        Para que se alcance el equilibrio, los productos de la reacción deben permanecer en el recipiente.
        En algunos casos se consigue sin necesidad de cerrarlo, por ejemplo, en las reacciones entre especies
        químicas solubles que no originen desprendimientos gaseosos. En cambio, en las reacciones
        entre gases o en las que se produzcan productos gaseosos, el recipiente deberá ser cerrado para
        evitar que escapen.
    \end{solu}
\end{exercise}

\begin{exercise} %?\problema{\marcador{Ejercicio 35, página 191}}
    Indica los factores que afecta al equilibrio de una reacción aplícalo al equilibrio de la síntesis de amoniaco, a partir de hidrógeno y nitrógeno.
    \begin{solu}
        El equilibrio de una reacción puede verse afectado por variaciones de la temperatura, de la presión
        o de la concentración de los productos y de los reactivos. Cuando esto ocurre, la reacción evoluciona
        en uno u otro sentido hasta alcanzar una nueva situación de equilibrio. El desplazamiento
        se produce precisamente en el sentido que contrarresta la perturbación (principio de Le Châtelier).
        \begin{reaction*}
            N2 \gas{} + 3 H2 \gas{} <-> 2 NH3 \gas{}
        \end{reaction*}
        La entalpía de esta reacción es: \( \Delta H = \qty{-92.4}{\kilo\joule\per\mole} \).

        De acuerdo con el principio de Le Châtelier, la formación de amoniaco se verá favorecida por una
        adición de reactivos (N2 y H2), por una retirada del producto (NH3), por un aumento de presión (hay
        menos moles gaseosos de productos que de reactivos) y por una disminución de la temperatura
        (al ser la reacción exotérmica). Sin embargo, si desciende demasiado la temperatura, el proceso es
        demasiado lento y ello obliga a trabajar a temperaturas
        de compromiso entre el factor termodinámico
        y el cinético.
    \end{solu}
\end{exercise}

\begin{exercise} %?\problema{\marcador{Ejercicio 44, página 191}}
    AB se disocia según: \formula{AB\ (g) <=> A\ (g) + B\ (g)} con \(\Delta H > 0\).
    ¿Qué sucede si aumenta la temperatura?, ¿y si disminuye el volumen?, ¿se alterará la constante de equilibrio?
    \begin{solu}
        Según el principio de Le Châtelier, las perturbaciones sobre el equilibrio desplazan a este en el
        sentido que tiende a corregirlas.
        \begin{itemize}
            \item Un aumento de temperatura desplazará el equilibrio en el sentido que tienda a rebajarla. Por tanto, hacia la derecha, en el sentido en el que la reacción es endotérmica (\( \Delta H > 0 \)).
            \item Si disminuye el volumen, el equilibrio se desplazará hacia la izquierda, ya que es menor el número
            de moles gaseosos en los reactivos (AB) que en los productos (A + B), según se aprecia en la estequiometría
            de la reacción.
            \item La variación de la presión o la disminución del volumen (apartado b) no alteran la constante de
            equilibrio. En cambio, si la reacción se lleva a cabo a otra temperatura (apartado a), sí cambia
            el valor de esta constante.
        \end{itemize}
    \end{solu}
\end{exercise}

\begin{exercise} %?\problema{\marcador{Ejercicio 46, página 191}}
    Si el siguiente equilibrio posee \(\Delta H = \energia[kj]{-100}\)
    \begin{center}
        \ch{SO3\ (g) + 1/2 O2\ (g) <=> SO3\ (g)}
    \end{center}
    predice en qué sentido se desplazará el equilibrio cuando se perturbe por:
    \begin{enumerate}[label=\letraminuscula]
        \item Aumento del volumen del recipiente.
        \item Adición de \ch{SO3}.
        \item Aumento de la temperatura.
        \item Adición de oxígeno.
        \item ¿En qué casos variará el valor de la constante de equilibrio?
    \end{enumerate}
    \begin{solu}
        \begin{enumerate}[label=\letraminuscula]
            \item Si aumenta el volumen, disminuye la presión. La reacción se desplazará en el sentido que tienda
            a contrarrestar
            dicha disminución. Concretamente, hacia la izquierda, produciendo un aumento
            del número de moles que aumente la presión.
            \item Un aumento de la concentración de SO3 desplazará la reacción hacia la izquierda, propiciando
            el consumo de SO3.
            \item Al ser la reacción exotérmica hacia la derecha (\(\Delta H < 0  \)), un aumento de temperatura la desplazará
            hacia la izquierda propiciando la absorción de calor.
            \item Si aumenta la [O2] se desplazará hacia la derecha, hacia el consumo del exceso de O2.
            \item Solo en el apartado c), puesto que la constante de equilibrio solo depende de la temperatura.
            47
        \end{enumerate}
    \end{solu}
\end{exercise}

\begin{exercise}    %?\problema{\marcador{Ejercicio 43, página 191}}
    Enuncia el principio de Le Châtelier.
    Si en una reacción: \formula{A + B <=> AB}, en fase gaseosa, la \(Kp\) vale \num{4,3} a \temperatura[c]{250} y tiene un valor de \num{1,8} acción será exotérmica o endotérmica y en qué sentido se desplazará el equilibrio al aumentar la temperatura
    \begin{solu}
%        \begin{figure}[H]
%            \centering
%            \includegraphics[scale=0.8]{images/imagesUD6EquilibrioQuimico/UD6_Problemas/Ejercicio_general_06}
%        \end{figure}
        El \textbf{principio de Le Châtelier} establece que si un sistema en equilibrio se ve sometido a una perturbación externa (cambio de concentración, presión, volumen o temperatura), el sistema evoluciona espontáneamente en el sentido que tiende a contrarrestar dicha perturbación, hasta alcanzar un nuevo estado de equilibrio.

        En este caso, la constante \(K_p\) de la reacción \(\ch{A + B <=> AB}\) \textbf{disminuye} al aumentar la temperatura (pasa de \(4{,}3\) a \(250\) ºC a \(1{,}8\) a una temperatura superior). Como un aumento de temperatura desplaza siempre el equilibrio en el sentido endotérmico, y aquí ese aumento de temperatura hace disminuir \(K_p\) (es decir, desplaza el equilibrio hacia la izquierda, hacia los reactivos A y B), se deduce que la reacción \textbf{directa} (formación de AB) es \textbf{exotérmica} (\(\Delta H<0\)), y la reacción inversa (disociación de AB) es endotérmica.

        Por tanto, al aumentar la temperatura, el equilibrio se desplaza hacia la \textbf{izquierda} (hacia A + B), favoreciendo el sentido endotérmico, en consonancia con la disminución observada de \(K_p\).
    \end{solu}
\end{exercise}

\begin{exercise}    %?\problema{\marcador{Ejercicio 49, página 192}}
    A 227 °C y 10 atm, el pentacloruro de fósforo, \ch{PCl5}, se disocia en tricloruro de fósforo, \ch{PCl3}, y cloro, \ch{Cl2}.
    La mezcla de reacción en equilibrio contiene un 20 \% en volumen de cloro.
    Calcula \(K_p\) y \(K_c\).
    \begin{solu}
%        \begin{figure}[H]
%            \centering
%            \includegraphics[scale=0.8]{images/imagesUD6EquilibrioQuimico/UD6_Problemas/Ejercicio_general_07}
%        \end{figure}
        La reacción de disociación es \(\ch{PCl5 (g) <=> PCl3 (g) + Cl2 (g)}\). Tomando como base 1 mol inicial de \ch{PCl5} y llamando \(\alpha\) al grado de disociación:

        \begin{center}
            \begin{tabular}{lccc}
                & \ch{PCl5} & \ch{PCl3} & \ch{Cl2} \\
                Equilibrio & \(1-\alpha\) & \(\alpha\) & \(\alpha\)
            \end{tabular}
        \end{center}

        Moles totales en el equilibrio: \(1+\alpha\). Como el \ch{Cl2} supone el 20 \% en volumen (fracción molar, ya que en una mezcla de gases ideales la fracción en volumen coincide con la fracción molar):
        \[
            x_{\ch{Cl2}} = \frac{\alpha}{1+\alpha} = 0{,}20 \quad \Rightarrow \quad \alpha = 0{,}20 + 0{,}20\,\alpha \quad \Rightarrow \quad \alpha = \frac{0{,}20}{0{,}80} = 0{,}25
        \]

        Moles totales: \(1+\alpha = 1{,}25\). Las fracciones molares son \(x_{\ch{PCl5}} = 0{,}75/1{,}25 = 0{,}6\), \(x_{\ch{PCl3}} = x_{\ch{Cl2}} = 0{,}25/1{,}25 = 0{,}2\).

        Con \(P_{total} = \presion{10}\) (10 atm), las presiones parciales son:
        \[
            P_{\ch{PCl5}} = 0{,}6 \times 10 = 6\ \text{atm}, \qquad P_{\ch{PCl3}} = P_{\ch{Cl2}} = 0{,}2 \times 10 = 2\ \text{atm}
        \]

        La constante \(K_p\) es:
        \[
            K_p = \frac{P_{\ch{PCl3}} \cdot P_{\ch{Cl2}}}{P_{\ch{PCl5}}} = \frac{2 \times 2}{6} = 0{,}667\ \text{atm}
        \]

        Para calcular \(K_c\), con \(\Delta n = 2-1 = 1\), \(T = 227+273 = 500\ \text{K}\) y \(R = \ctegases{}\):
        \[
            K_c = \frac{K_p}{(RT)^{\Delta n}} = \frac{0{,}667}{(0{,}082 \times 500)^1} = \frac{0{,}667}{41} = \num{0.0163}\ \text{mol/L}
        \]
    \end{solu}
\end{exercise}

\begin{exercise}    %?\problema{\marcador{Ejercicio 50, página 192}}
    A 100 °C y 1 atm, el hierro es oxidado por el vapor de agua hasta óxido ferroso y se desprende hidrógeno.
    Cuando se llega al equilibrio en la reacción, la presión parcial del hidrógeno es de 490 mmHg, y la del agua, 270 mmHg. Si el
    grado de disociación del agua a \temperatura[c]{100} y a la presión de 1 atm vale \num{6.4e-7}, calcula la presión de disociación del
    óxido ferroso.
    \begin{solu}
%        \begin{figure}[H]
%            \centering
%            \includegraphics[scale=0.8]{images/imagesUD6EquilibrioQuimico/UD6_Problemas/Ejercicio_general_08}
%        \end{figure}
        Se combinan dos equilibrios mediante la ley de Hess (los equilibrios se pueden sumar o restar igual que las reacciones termoquímicas, multiplicando o dividiendo sus constantes):

        \textbf{Equilibrio A} (oxidación del hierro): \(\ch{Fe (s) + H2O (g) <=> FeO (s) + H2 (g)}\)
        \[
            K_A = \frac{P_{\ch{H2}}}{P_{\ch{H2O}}} = \frac{490}{270} = 1{,}815
        \]

        \textbf{Equilibrio B} (disociación del agua): \(\ch{H2O (g) <=> H2 (g) + 1/2 O2 (g)}\), con grado de disociación \(\alpha = \num{6.4e-7}\) a 1 atm. Tomando 1 mol inicial de \ch{H2O}, en el equilibrio (con \(\alpha\) tan pequeño, los moles totales \(\approx 1\)):
        \[
            P_{\ch{H2O}} \approx (1-\alpha)\cdot 1 \approx 1\ \text{atm}, \quad P_{\ch{H2}} = \alpha \cdot 1 = \num{6.4e-7}\ \text{atm}, \quad P_{\ch{O2}} = \frac{\alpha}{2}\cdot 1 = \num{3.2e-7}\ \text{atm}
        \]
        \[
            K_B = \frac{P_{\ch{H2}} \cdot P_{\ch{O2}}^{1/2}}{P_{\ch{H2O}}} = \frac{\num{6.4e-7} \times \sqrt{\num{3.2e-7}}}{1} = \num{3.62e-10}
        \]

        La presión de disociación del óxido ferroso corresponde al equilibrio \textbf{C}: \(\ch{FeO (s) <=> Fe (s) + 1/2 O2 (g)}\), con \(K_C = P_{\ch{O2}}^{1/2}\). Esta reacción se obtiene como (B) \(-\) (A) (equilibrio B menos equilibrio A, ya que al sumar B con el inverso de A se cancelan \ch{H2O} y \ch{H2}):
        \[
            K_C = \frac{K_B}{K_A} = \frac{\num{3.62e-10}}{1{,}815} = \num{1.995e-10}
        \]

        Por tanto, la presión de disociación del óxido ferroso (presión parcial de \ch{O2} en el equilibrio \ch{FeO <=> Fe + 1/2 O2}) es:
        \[
            P_{\ch{O2}} = K_C^2 = \left(\num{1.995e-10}\right)^2 = \num{3.98e-20}\ \text{atm}
        \]

        Este valor tan extremadamente pequeño indica que el óxido ferroso es una sustancia muy estable, que prácticamente no se descompone en Fe y O2 en condiciones normales.
    \end{solu}
\end{exercise}

\begin{exercise}
    Deduce si se formará precipitado de cloruro de plata cuyo \(K_s = \num{1.7e-10}\) a \temperatura[c]{25} al añadir a \volumen[centimetroscubicos]{250} de cloruro de sodio \concentracion[m]{0.02} \volumen[centimetroscubicos]{50} de nitrato de plata \concentracion[m]{0.5}
    \begin{solu}

        Al mezclar ambas disoluciones, los moles de cada ion no varían (no hay reacción entre ellos hasta comprobar si precipita), pero sí su concentración, al diluirse en el volumen total de la mezcla.

        Moles de \ch{Cl-} aportados por el \ch{NaCl}:

        \begin{equation*}
            n(\ch{Cl-}) = 0,250 \ \si{\litre} \times 0,02 \ \si{\Molar} = \num{5e-3} \ \text{mol}
        \end{equation*}

        Moles de \ch{Ag+} aportados por el \ch{AgNO3}:

        \begin{equation*}
            n(\ch{Ag+}) = 0,050 \ \si{\litre} \times 0,5 \ \si{\Molar} = \num{2.5e-2} \ \text{mol}
        \end{equation*}

        El volumen total de la mezcla es \(V = 250 + 50 = \SI{300}{\centi\metre\cubed} = \SI{0.300}{\litre}\), por lo que las concentraciones tras la mezcla son:

        \begin{equation*}
            [\ch{Cl-}] = \dfrac{5 \times 10^{-3}}{0,300} = \num{1.67e-2} \ \si{\Molar} \qquad [\ch{Ag+}] = \dfrac{2,5 \times 10^{-2}}{0,300} = \num{8.33e-2} \ \si{\Molar}
        \end{equation*}

        Calculamos el producto iónico \(Q\) con estas concentraciones y lo comparamos con \(K_s\):

        \begin{equation*}
            Q = [\ch{Ag+}][\ch{Cl-}] = (8,33 \times 10^{-2})(1,67 \times 10^{-2}) = 1,39 \times 10^{-3}
        \end{equation*}

        Como \(Q = 1,39 \times 10^{-3} \gg K_s = 1,7 \times 10^{-10}\), la disolución está sobresaturada respecto al \ch{AgCl}, por lo que \textbf{sí se formará precipitado} de cloruro de plata.

    \end{solu}
\end{exercise}

\subsection{Ejercicios PAU}
\begin{exerciseEVAU}
    Conociendo la solubilidad del \ch{Ca(IO3)2}, calcule su producto de solubilidad.

    Datos: s = \qty{5,6e-3}{\M}
\begin{solu}

\end{solu}
\end{exerciseEVAU}
\begin{exerciseEVAU}
    El fosgeno (\ch{COCl2}) es un producto gaseoso que se descompone en monóxido de
carbono y cloro según el proceso:
    \begin{equation*}
        \ch{COCl2 \gas{} <=> CO \gas{} + Cl2 \gas{}}
    \end{equation*}
En un recipiente de 250 mL de capacidad se introdujeron 0,213 gramos de fosgeno, de manera que cuando
se alcanzó el equilibrio a la temperatura de 27 ºC, la presión en el interior del matraz fue de 230 mm de
Hg. A partir de estos datos, calcule:
    \begin{enumerate}[label=\letraminuscula]
        \item El grado de disociación del fosgeno.
        \item La presión parcial de cada componente gaseoso en la mezcla.
        \item El valor de las constantes Kp y Kc.
    \end{enumerate}

\begin{solu}

\end{solu}
\end{exerciseEVAU}

\begin{exerciseEVAU}
    En un recipiente cerrado y vacío de 5 L de capacidad, a 727ºC, se introducen 1 mol de selenio y 1 mol de
hidrógeno, alcanzándose el equilibrio siguiente:
    \begin{centeredequation}
        Se \gas{} + H2 \gas{} <-> H2Se \gas{}
    \end{centeredequation}
    Cuando se alcanza el equilibrio se observa que la presión en el interior del recipiente es de 18,1 atm.
    \begin{enumerate}[label=\letraminuscula]
        \item Calcule las concentraciones de cada una de las especies en el equilibrio.
        \item Calcule el valor de Kp y de Kc.
    \end{enumerate}
\begin{solu}

    La reacción es \(\ch{Se \gas{} + H2 \gas{} <=> H2Se \gas{}}\), cuyo número de moles gaseosos totales disminuye en 1 por cada mol de \ch{H2Se} formado (2 mol de reactivos por 1 mol de producto).

    \textbf{a) Concentraciones en el equilibrio}

    Con \(T = \SI{1000}{\kelvin}\) (\(727+273\)), calculamos los moles totales en el equilibrio a partir de la presión total:

    \begin{equation*}
        n_T = \dfrac{P_T\, V}{R\, T} = \dfrac{18,1 \times 5}{0,082 \times 1000} = 1,104 \ \text{mol}
    \end{equation*}

    Si llamamos \(x\) a los moles de \ch{Se} (y de \ch{H2}) que reaccionan, en el equilibrio: \(n(\ch{Se}) = 1-x\), \(n(\ch{H2}) = 1-x\), \(n(\ch{H2Se}) = x\), con \(n_T = (1-x)+(1-x)+x = 2-x\).

    \begin{equation*}
        2 - x = 1,104 \ \Rightarrow \ x = 0,896 \ \text{mol}
    \end{equation*}

    Por tanto: \(n(\ch{Se}) = n(\ch{H2}) = 1 - 0,896 = 0,104\) mol, y \(n(\ch{H2Se}) = 0,896\) mol. Dividiendo entre el volumen (5 L):

    \begin{equation*}
        [\ch{Se}] = [\ch{H2}] = \dfrac{0,104}{5} = \SI{0.0208}{\Molar} \qquad [\ch{H2Se}] = \dfrac{0,896}{5} = \SI{0.1793}{\Molar}
    \end{equation*}

    \textbf{b) Kp y Kc}

    \begin{equation*}
        K_c = \dfrac{[\ch{H2Se}]}{[\ch{Se}][\ch{H2}]} = \dfrac{0,1793}{(0,0208)^2} = 414{,}3
    \end{equation*}

    Como \(\Delta n = 1 - 2 = -1\):

    \begin{equation*}
        K_p = K_c\, (R T)^{\Delta n} = \dfrac{K_c}{R T} = \dfrac{414,3}{0,082 \times 1000} = 5,05
    \end{equation*}

\end{solu}
\end{exerciseEVAU}

\begin{exerciseEVAU}
    En un matraz de un litro se introducen 0,2 moles de \ch{PCl5} y se calienta a 300ºC. A esta temperatura, el \ch{PCl5} se disocia según la reacción mostrada, y la presión parcial de \ch{PCl3} observada en equilibrio fue de 0,5 atm.
    \begin{centeredequation}
        \ch{PCl5 <-> PCl3 + Cl2}
    \end{centeredequation}
    \begin{enumerate}[label=\letraminuscula]
    \item Calcule el número de moles de todas las especies en el equilibrio, y el valor de la presión
    parcial del \ch{PCl5} en equilibrio.
    \item Calcule el grado de disociación del \ch{PCl5}.
    \item Calcule Kp y Kc a 300ºC
    \end{enumerate}
    \begin{solu}

        La reacción es \(\ch{PCl5 \gas{} <=> PCl3 \gas{} + Cl2 \gas{}}\), con \(T = \SI{573}{\kelvin}\) (\(300+273\)) y \(V = \SI{1}{\litre}\).

        \textbf{a) Moles en el equilibrio y presión parcial del \ch{PCl5}}

        Si llamamos \(x\) a los moles de \ch{PCl5} que se disocian, \(n(\ch{PCl3}) = n(\ch{Cl2}) = x\). Usando la presión parcial de \ch{PCl3} (gas ideal, \(V=1\) L):

        \begin{equation*}
            x = n(\ch{PCl3}) = \dfrac{p_{\ch{PCl3}}\, V}{R\, T} = \dfrac{0,5 \times 1}{0,082 \times 573} = 0,01064 \ \text{mol}
        \end{equation*}

        Por tanto \(n(\ch{PCl5})_{eq} = 0,2 - 0,01064 = 0,18936\) mol, y su presión parcial:

        \begin{equation*}
            p_{\ch{PCl5}} = \dfrac{n(\ch{PCl5})\, R\, T}{V} = \dfrac{0,18936 \times 0,082 \times 573}{1} = \SI{8.90}{\atmospher}
        \end{equation*}

        \textbf{b) Grado de disociación}

        \begin{equation*}
            \alpha = \dfrac{x}{n_0} = \dfrac{0,01064}{0,2} = 0,0532 \ (5,32\ \%)
        \end{equation*}

        \textbf{c) Kp y Kc}

        Con \(V=1\) L, las concentraciones coinciden numéricamente con los moles: \([\ch{PCl5}] = 0,18936\ \si{\Molar}\), \([\ch{PCl3}] = [\ch{Cl2}] = 0,01064\ \si{\Molar}\).

        \begin{equation*}
            K_c = \dfrac{[\ch{PCl3}][\ch{Cl2}]}{[\ch{PCl5}]} = \dfrac{(0,01064)^2}{0,18936} = \num{5.98e-4} \ \si{\Molar}
        \end{equation*}

        Como \(\Delta n = 2 - 1 = 1\):

        \begin{equation*}
            K_p = K_c\, (R T) = 5,98 \times 10^{-4} \times (0,082 \times 573) = 0,0281 \ \si{\atmospher}
        \end{equation*}

        (Puede comprobarse directamente: \(K_p = p_{\ch{PCl3}}\, p_{\ch{Cl2}} / p_{\ch{PCl5}} = 0,5^2/8,90 = 0,0281\), coherente.)

    \end{solu}
\end{exerciseEVAU}

\begin{exerciseEVAU}
    En un laboratorio de control de calidad de aguas, se emplea el equilibrio de precipitación del cromato de
    plata (\ch{Ag2CrO4}) para detectar la presencia de iones plata en muestras contaminadas. Para ello, se trabaja
    con disoluciones que ya contienen cromato potásico (\ch{K2CrO4}) como reactivo, a una concentración de 0,01
    M.
    \begin{enumerate}[label=\letraminuscula]
    \item Escriba el equilibrio de disolución del \ch{Ag2CrO4} en agua pura, la expresión de su producto de
    solubilidad y la solubilidad de \ch{Ag2CrO4} en mol/L.
    \item  La solubilidad del \ch{Ag2CrO4} en presencia de la sal de \ch{K2CrO4}, expresada en mol/L y en g/L.
    \end{enumerate}
    \textbf{Datos:} \( K_s \ch{Ag2CrO2} = \num{3.9e-12} \)
    \begin{solu}

        \textbf{a) Equilibrio, producto de solubilidad y solubilidad en agua pura}

        \begin{equation*}
            \ch{Ag2CrO4 \sld{} <=> 2 Ag+ \aq{} + CrO4^2- \aq{}} \qquad K_s = [\ch{Ag+}]^2 [\ch{CrO4^2-}]
        \end{equation*}

        Si \(s\) es la solubilidad, \([\ch{Ag+}] = 2s\) y \([\ch{CrO4^2-}] = s\), por lo que:

        \begin{equation*}
            K_s = (2s)^2 (s) = 4 s^3 = 3,9 \times 10^{-12}
        \end{equation*}

        \begin{equation*}
            s = \sqrt[3]{\dfrac{3,9 \times 10^{-12}}{4}} = \sqrt[3]{9,75 \times 10^{-13}} = \num{9.92e-5} \ \si{\Molar}
        \end{equation*}

        \textbf{b) Solubilidad en presencia de \ch{K2CrO4} 0,01 M}

        El \ch{K2CrO4} aporta un ion común, \ch{CrO4^2-}. Si \(s'\) es la nueva solubilidad, \([\ch{Ag+}] = 2s'\) y \([\ch{CrO4^2-}] = 0,01 + s' \approx 0,01\ \si{\Molar}\) (ya que, por el efecto del ion común, \(s'\) será mucho menor que 0,01):

        \begin{equation*}
            K_s = (2s')^2 (0,01) = 3,9 \times 10^{-12}
        \end{equation*}

        \begin{equation*}
            s'^2 = \dfrac{3,9 \times 10^{-12}}{4 \times 0,01} = 9,75 \times 10^{-11} \ \Rightarrow \ s' = \num{9.87e-6} \ \si{\Molar}
        \end{equation*}

        Comprobamos que la aproximación es válida, ya que \(s' = 9,87 \times 10^{-6} \ll 0,01\).

        Expresando la solubilidad en g/L, con \(M(\ch{Ag2CrO4}) = 2(107,87) + 52,00 + 4(16,00) = \SI{331.74}{\gram\per\mole}\):

        \begin{equation*}
            s' = 9,87 \times 10^{-6} \ \text{mol/L} \times 331,74 \ \si{\gram\per\mole} = \num{3.27e-3} \ \si{\gram\per\litre}
        \end{equation*}

        Como cabía esperar por el efecto del ion común, la solubilidad del \ch{Ag2CrO4} en presencia de \ch{K2CrO4} (\(9,87 \times 10^{-6}\) M) es muy inferior a la que tiene en agua pura (\(9,92 \times 10^{-5}\) M).

    \end{solu}
\end{exerciseEVAU}

\begin{exerciseEVAU}

    La empresa AgroFer S.A. busca optimizar la producción de amoníaco (\ch{NH3}) mediante el proceso HaberBosch, clave en la fabricación de fertilizantes. Dado su alto consumo energético, pequeñas mejoras pueden tener gran impacto económico y ambiental. Este proceso combina nitrógeno del aire (\ch{N2}) y gas hidrógeno (\ch{H2}) a alta presión y temperatura, dando lugar a amoníaco (\ch{NH3}), un compuesto fundamental para la producción de fertilizantes que permiten mejorar los rendimientos agrícolas. Como parte del equipo técnico, ¿podría analizar cómo afectan distintas condiciones al rendimiento de la síntesis de amoníaco (\ch{NH3}) para proponer soluciones eficientes y sostenibles?. En la planta piloto se introduce una mezcla de gases en un reactor de 10 L a 400 °C con las cantidades de 1 mol de nitrógeno (\ch{N2}) y 3 moles de hidrógeno (\ch{H2}) como cantidades iniciales. Cuando se alcanza el equilibrio, se ha formado 1 mol de \ch{NH3} (g) y se ha observado que la entalpía obtenida (\( \Delta H \)) en la reacción es de -92 KJ/mol.
    Las cuestiones a resolver son las siguientes:
    \begin{enumerate}[label=\letraminuscula]
        \item Formular la reacción de síntesis del amoníaco según los datos que indica el enunciado
        y ajustarla.
        \item Determinar el valor de las constantes de equilibrio Kc y Kp a 400 °C.
        \item Explicar cómo se vería afectado el equilibrio si se aumentara la presión del sistema,
        manteniendo la temperatura constante. Y analizar qué ocurriría si se eleva la temperatura del
        proceso.
    \end{enumerate}

    \begin{solu}

        \textbf{a) Reacción ajustada}

        \begin{equation*}
            \ch{N2 \gas{} + 3 H2 \gas{} <=> 2 NH3 \gas{}}
        \end{equation*}

        \textbf{b) Constantes de equilibrio a 400 ºC}

        Con \(T = \SI{673}{\kelvin}\) (\(400+273\)) y \(V = \SI{10}{\litre}\). Como se forma 1 mol de \ch{NH3} y la estequiometría indica que por cada mol de \ch{N2} que reacciona se forman 2 mol de \ch{NH3}, el número de moles de \ch{N2} que reaccionan es \(x = 1/2 = 0,5\) mol.

        \begin{equation*}
            n(\ch{N2})_{eq} = 1 - 0,5 = 0,5 \ \text{mol} \qquad n(\ch{H2})_{eq} = 3 - 3(0,5) = 1,5 \ \text{mol} \qquad n(\ch{NH3})_{eq} = 1 \ \text{mol}
        \end{equation*}

        Las concentraciones en el equilibrio son:

        \begin{equation*}
            [\ch{N2}] = \dfrac{0,5}{10} = \SI{0.05}{\Molar} \qquad [\ch{H2}] = \dfrac{1,5}{10} = \SI{0.15}{\Molar} \qquad [\ch{NH3}] = \dfrac{1}{10} = \SI{0.1}{\Molar}
        \end{equation*}

        \begin{equation*}
            K_c = \dfrac{[\ch{NH3}]^2}{[\ch{N2}][\ch{H2}]^3} = \dfrac{(0,1)^2}{0,05 \times (0,15)^3} = 59,3 \ \si{\square\litre\per\square\mole}
        \end{equation*}

        Como \(\Delta n = 2 - 4 = -2\):

        \begin{equation*}
            K_p = K_c\, (R T)^{-2} = \dfrac{59,3}{(0,082 \times 673)^2} = \dfrac{59,3}{3045,9} = 0,0195 \ \si{\per\square\atmospher}
        \end{equation*}

        \textbf{c) Efecto de la presión y de la temperatura}

        Al \textbf{aumentar la presión} manteniendo constante la temperatura, el equilibrio se desplaza en el sentido en el que hay menos moles de gas, es decir, \textbf{hacia la derecha} (2 mol de \ch{NH3} frente a 4 mol de \ch{N2}+\ch{H2}), aumentando el rendimiento de la síntesis de amoníaco.

        Como la reacción es exotérmica (\(\Delta H = \energia[kjmol]{-92}\)), al \textbf{elevar la temperatura} el equilibrio se desplaza hacia la izquierda (sentido endotérmico), \textbf{disminuyendo} la cantidad de \ch{NH3} formada, aunque la reacción alcanza el equilibrio más rápidamente. Por ello, en la práctica industrial (proceso Haber-Bosch) se trabaja a una temperatura de compromiso, ni demasiado baja (poca velocidad) ni demasiado alta (poco rendimiento), junto con un catalizador para acelerar el proceso sin necesidad de subir la temperatura.

    \end{solu}

\end{exerciseEVAU}

\begin{exerciseEVAU}
    \textbf{2024}. A una temperatura de \SI{460}{\celsius}  se introdujeron \num{2,5} moles de \ch{NOCl} en un recipiente cerrado de volumen \SI{1}{\litre}.
    Una vez alcanzado el equilibrio se determinó que se habían formado \num{0,78} moles de \ch{NO}.
    \begin{reaction*}
        2 NOCl \gas{} <=> Cl2 \gas{} + 2 NO \gas{}
    \end{reaction*}
    \begin{enumerate}[label=\letraminuscula]
        \item Calcule la constante de equilibrio Kc.
        \item Calcule la constante de equilibrio Kp.
        \item Calcule la presión total de la mezcla en equilibrio a esa temperatura.
    \end{enumerate}

    \datos{}: \ctegases{}

    \begin{solu}

        La reacción es \(\ch{2 NOCl \gas{} <=> Cl2 \gas{} + 2 NO \gas{}}\), con \(T = \SI{733}{\kelvin}\) (\(460+273\)) y \(V = \SI{1}{\litre}\).

        Como se han formado 0,78 mol de \ch{NO}, y por estequiometría \(n(\ch{NO}) = 2x\) siendo \(x\) el número de moles de \ch{Cl2} formados:

        \begin{equation*}
            x = \dfrac{0,78}{2} = 0,39 \ \text{mol}
        \end{equation*}

        \begin{equation*}
            n(\ch{NOCl})_{eq} = 2,5 - 2(0,39) = 1,72 \ \text{mol} \qquad n(\ch{Cl2})_{eq} = 0,39 \ \text{mol} \qquad n(\ch{NO})_{eq} = 0,78 \ \text{mol}
        \end{equation*}

        Como \(V = 1\) L, las concentraciones coinciden numéricamente con los moles.

        \textbf{a) Constante \(K_c\)}

        \begin{equation*}
            K_c = \dfrac{[\ch{Cl2}][\ch{NO}]^2}{[\ch{NOCl}]^2} = \dfrac{0,39 \times (0,78)^2}{(1,72)^2} = 0,0802 \ \si{\Molar}
        \end{equation*}

        \textbf{b) Constante \(K_p\)}

        Como \(\Delta n = 3 - 2 = 1\):

        \begin{equation*}
            K_p = K_c\, (R T) = 0,0802 \times (0,082 \times 733) = 4,82 \ \si{\atmospher}
        \end{equation*}

        \textbf{c) Presión total en el equilibrio}

        \begin{equation*}
            n_T = 1,72 + 0,39 + 0,78 = 2,89 \ \text{mol}
        \end{equation*}

        \begin{equation*}
            P_T = \dfrac{n_T\, R\, T}{V} = \dfrac{2,89 \times 0,082 \times 733}{1} = \SI{173.7}{\atmospher}
        \end{equation*}

    \end{solu}

\end{exerciseEVAU}



\begin{exerciseEVAU}
    \textbf{2024}. Sabiendo que el producto de solubilidad de \ch{CaF2} es \( \Kps = \num{3.4e-11} \) a \SI{25}{\celsius}, calcule la solubilidad molar de \ch{CaF2} a esa temperatura.

    \begin{solu}

        El equilibrio de solubilidad es:

        \begin{equation*}
            \ch{CaF2 \sld{} <=> Ca^2+ \aq{} + 2 F- \aq{}}
        \end{equation*}

        Si \(s\) es la solubilidad molar, en el equilibrio \([\ch{Ca^2+}] = s\) y \([\ch{F-}] = 2s\), por lo que:

        \begin{equation*}
            K_{ps} = [\ch{Ca^2+}][\ch{F-}]^2 = s\, (2s)^2 = 4 s^3
        \end{equation*}

        \begin{equation*}
            s = \sqrt[3]{\dfrac{K_{ps}}{4}} = \sqrt[3]{\dfrac{3,4 \times 10^{-11}}{4}} = \sqrt[3]{8,5 \times 10^{-12}} = \num{2.04e-4} \ \si{\Molar}
        \end{equation*}

    \end{solu}
\end{exerciseEVAU}

\begin{exerciseEVAU}
    \textbf{2023}.
    En un recipiente de 2L se introducen \SI{92,4}{\g} de \ch{CO2} y \SI{3,2}{\g} de \ch{H2}, calentándose la mezcla a 1800 ºC. Una vez alcanzado el siguiente equilibrio: \ch{CO2 \gas{} + H2 \gas{} <=> CO \gas{} + H2O \gas{}} se analiza la mezcla, encontrándose que quedan \num{0,9} moles de \ch{CO2}.

    \begin{enumerate}[label=\letraminuscula]
        \item Calcule la concentración de cada especie en el equilibrio.
        \item Calcule \Kc{} y \Kp{} a \SI{1800}{\celsius}.
        \item Justifique si una disminución del volumen del recipiente afectaría al equilibrio.
    \end{enumerate}

    \datos{}: Masas atómicas (u): C = 12, O = 16, H = 1; \(R = \ctegases\)

    \begin{solu}

        La reacción es \(\ch{CO2 \gas{} + H2 \gas{} <=> CO \gas{} + H2O \gas{}}\), con \(T = \SI{2073}{\kelvin}\) (\(1800+273\)) y \(V = \SI{2}{\litre}\).

        Los moles iniciales son:

        \begin{equation*}
            n_0(\ch{CO2}) = \dfrac{92,4}{44} = 2,1 \ \text{mol} \qquad n_0(\ch{H2}) = \dfrac{3,2}{2} = 1,6 \ \text{mol}
        \end{equation*}

        Como quedan 0,9 mol de \ch{CO2} en el equilibrio, han reaccionado \(x = 2,1 - 0,9 = 1,2\) mol (estequiometría 1:1:1:1):

        \begin{equation*}
            n(\ch{H2})_{eq} = 1,6 - 1,2 = 0,4 \ \text{mol} \qquad n(\ch{CO})_{eq} = n(\ch{H2O})_{eq} = 1,2 \ \text{mol}
        \end{equation*}

        \textbf{a) Concentraciones en el equilibrio}

        \begin{equation*}
            [\ch{CO2}] = \dfrac{0,9}{2} = \SI{0.45}{\Molar} \qquad [\ch{H2}] = \dfrac{0,4}{2} = \SI{0.2}{\Molar} \qquad [\ch{CO}] = [\ch{H2O}] = \dfrac{1,2}{2} = \SI{0.6}{\Molar}
        \end{equation*}

        \textbf{b) \(K_c\) y \(K_p\) a 1800 ºC}

        \begin{equation*}
            K_c = \dfrac{[\ch{CO}][\ch{H2O}]}{[\ch{CO2}][\ch{H2}]} = \dfrac{0,6 \times 0,6}{0,45 \times 0,2} = 4,0
        \end{equation*}

        Como \(\Delta n = 2 - 2 = 0\), \(K_p = K_c\, (RT)^0 = K_c = 4,0\) (adimensional).

        \textbf{c) Efecto de una disminución del volumen}

        Como el número de moles gaseosos es el mismo en reactivos que en productos (\(\Delta n = 0\)), una variación de la presión (por ejemplo, al disminuir el volumen del recipiente) \textbf{no desplaza el equilibrio}, ya que ambos sentidos de la reacción se ven afectados por igual.

    \end{solu}
\end{exerciseEVAU}


\begin{exerciseEVAU}
    \textbf{2023} Sabiendo que a \SI{298}{\kelvin}, la solubilidad del \ch{CaBr2} en agua es \SI{2}{\mole\per\litre}:
    \begin{enumerate}[label=\letraminuscula]
        \item Calcule la constante del producto de solubilidad (\Kps{}) de dicha sal.
        \item Calcule la solubilidad del CaBr2 en presencia de una disolución \SI{0.2}{\M} de \ch{NaBr}.
    \end{enumerate}

    \begin{solu}

        \textbf{a) Producto de solubilidad}

        \begin{equation*}
            \ch{CaBr2 \sld{} <=> Ca^2+ \aq{} + 2 Br- \aq{}}
        \end{equation*}

        Con \(s = \SI{2}{\mole\per\litre}\), \([\ch{Ca^2+}] = s\) y \([\ch{Br-}] = 2s\):

        \begin{equation*}
            K_{ps} = [\ch{Ca^2+}][\ch{Br-}]^2 = s\, (2s)^2 = 4 s^3 = 4 (2)^3 = 32
        \end{equation*}

        \textbf{b) Solubilidad en presencia de \ch{NaBr} 0,2 M}

        Llamando \(s'\) a la nueva solubilidad, \([\ch{Ca^2+}] = s'\) y \([\ch{Br-}] = 0,2 + 2s'\) (el ion \ch{Br-} procede tanto del \ch{NaBr} como del \ch{CaBr2} disuelto):

        \begin{equation*}
            K_{ps} = s'\, (0,2 + 2s')^2 = 32
        \end{equation*}

        Dado que \(K_{ps}\) es muy grande, no puede despreciarse \(2s'\) frente a 0,2 M (la sal es muy soluble). Resolviendo la ecuación de forma numérica se obtiene:

        \begin{equation*}
            s' \approx \SI{1.93}{\mole\per\litre}
        \end{equation*}

        (Comprobación: \(1,93 \times (0,2 + 3,86)^2 = 1,93 \times 4,06^2 = 1,93 \times 16,48 \approx 31,8 \approx 32\).)

        Como cabe esperar por el efecto del ion común, la solubilidad en presencia de \ch{NaBr} (1,93 M) es ligeramente inferior a la solubilidad en agua pura (2 M).

    \end{solu}
\end{exerciseEVAU}

\begin{exerciseEVAU}
    \textbf{2023}. A una temperatura determinada, la \Kc{} para la descomposición del fluoruro de hidrógeno según la reacción:
    \begin{reaction*}
        2 HF \gas{} <-> H2 \gas{} + F2 \gas{}
    \end{reaction*}
    tiene un valor de \num{1.0e-13}.
    Al cabo de cierto tiempo se encuentran las siguientes concentraciones: [\ch{HF}] = \qty{0.5}{\M}; [\ch{H2}] = \qty{e-3}{\M}1; [\ch{F2}] = \qty{4e-3}{\M}.
    \begin{enumerate}[label=\letraminuscula]
        \item Justifique porqué el sistema no se encuentra en equilibrio.
        \item ¿Cómo ha de evolucionar el sistema para alcanzar el equilibrio
    \end{enumerate}

    \begin{solu}

        \textbf{a) El sistema no está en equilibrio}

        Para comprobar si el sistema está en equilibrio, calculamos el cociente de reacción \(Q\) con las concentraciones dadas y lo comparamos con \(K_c\):

        \begin{equation*}
            Q = \dfrac{[\ch{H2}][\ch{F2}]}{[\ch{HF}]^2} = \dfrac{(10^{-3})(4 \times 10^{-3})}{(0,5)^2} = \dfrac{4 \times 10^{-6}}{0,25} = 1,6 \times 10^{-5}
        \end{equation*}

        Como \(Q = 1,6 \times 10^{-5} \neq K_c = 1,0 \times 10^{-13}\), el sistema \textbf{no se encuentra en equilibrio}.

        \textbf{b) Evolución hacia el equilibrio}

        Como \(Q > K_c\) (de hecho, \(Q\) es unos 8 órdenes de magnitud mayor que \(K_c\)), hay un exceso de productos (\ch{H2} y \ch{F2}) respecto a la situación de equilibrio. Para que el sistema evolucione hacia el equilibrio, la reacción debe transcurrir en el sentido en el que \(Q\) disminuya hasta igualar a \(K_c\), es decir, \textbf{hacia la izquierda} (formación de \ch{HF} a partir de \ch{H2} y \ch{F2}), consumiendo \ch{H2} y \ch{F2} y aumentando \([\ch{HF}]\).

    \end{solu}
\end{exerciseEVAU}

\begin{exerciseEVAU}
    \textbf{2023}. En un recipiente cerrado de \qty{400}{\mL} 400 mL, se añaden \qty{1.280}{\g} de bromo (\ch{Br2}) y \qty{2.032}{\g} de yodo (\ch{I2}). Se eleva la temperatura a \qty{150}{\celsius}  y se alcanza el siguiente equilibrio:
    \begin{reaction*}
        Br2 \gas{} + I2 \gas{} -> 2 BrI \gas{}
    \end{reaction*}
    En estas condiciones el valor de \Kc es 280.
    \begin{enumerate}[label=\letraminuscula]
        \item Calcule el valor del grado de disociación (expresado en \%)
        \item Calcule el valor de Kp para este equilibrio a \qty{150}{\celsius}.
        \item Calcule los gramos de yodo en el equilibrio.
    \end{enumerate}

    \datos{}: Masas atómicas (u): Br = 80, I = 127; R = \ctegases{}

    \begin{solu}

        La reacción es \(\ch{Br2 \gas{} + I2 \gas{} <=> 2 BrI \gas{}}\), con \(T = \SI{423}{\kelvin}\) (\(150+273\)) y \(V = \SI{0.4}{\litre}\).

        Los moles iniciales son:

        \begin{equation*}
            n_0(\ch{Br2}) = \dfrac{1,280}{160} = 8 \times 10^{-3} \ \text{mol} \qquad n_0(\ch{I2}) = \dfrac{2,032}{254} = 8 \times 10^{-3} \ \text{mol}
        \end{equation*}

        Las concentraciones iniciales, iguales para ambos reactivos, son \([\ch{Br2}]_0 = [\ch{I2}]_0 = 8\times 10^{-3}/0,4 = \SI{0.02}{\Molar}\).

        \textbf{a) Grado de disociación}

        Llamando \(x\) a la concentración de \ch{Br2} (e \ch{I2}) que reacciona:

        \begin{equation*}
            K_c = \dfrac{(2x)^2}{(0,02-x)^2} = 280 \ \Rightarrow \ \dfrac{2x}{0,02-x} = \sqrt{280} = 16,73
        \end{equation*}

        \begin{equation*}
            2x = 16,73(0,02 - x) \ \Rightarrow \ 18,73\, x = 0,3347 \ \Rightarrow \ x = \num{1.787e-2} \ \si{\Molar}
        \end{equation*}

        \begin{equation*}
            \alpha = \dfrac{x}{0,02} = \dfrac{0,01787}{0,02} = 0,893 \ (89,3\ \%)
        \end{equation*}

        \textbf{b) Constante \(K_p\)}

        Como \(\Delta n = 2 - 2 = 0\), \(K_p = K_c\, (RT)^0 = K_c = 280\) (adimensional).

        \textbf{c) Gramos de yodo en el equilibrio}

        \begin{equation*}
            [\ch{I2}]_{eq} = 0,02 - 0,01787 = \num{2.13e-3} \ \si{\Molar}
        \end{equation*}

        \begin{equation*}
            n(\ch{I2})_{eq} = 2,13 \times 10^{-3} \times 0,4 = \num{8.54e-4} \ \text{mol}
        \end{equation*}

        \begin{equation*}
            m(\ch{I2}) = 8,54 \times 10^{-4} \times 254 = \SI{0.217}{\gram}
        \end{equation*}

    \end{solu}
\end{exerciseEVAU}

\begin{exerciseEVAU}
    En un recipiente de 1 litro de capacidad, en el que previamente se ha hecho el vacío, se introducen 0,1 mol de \ch{NO}; \moles{0,05} de \ch{H2} y \moles{0,1} de agua.
    Se calienta el matraz y se establece el equilibrio:
    \begin{center}
        \formula{2 NO\ (g) + 2 H2\ (g) <=> N2\ (g) + 2 H2O\ (g)}
    \end{center}
    Sabiendo que cuando se establece el equilibrio la concentración de NO es 0,062 M, calcule:
    \begin{enumerate}[label=\letraminuscula]
        \item Las concentraciones de todas las especies en el equilibrio.
        \item El valor de la constante \(K_c\) a esa temperatura.
    \end{enumerate}

% \solution{a)\(\conc{NO} = \concentracion[m]{0,062}\); \(\conc{H2} = \concentracion[m]{0,012}\); \(\conc{N2} = \concentracion[m]{0,019}\); \(\conc{H20} = \concentracion[m]{0,138}\). b) \(k_c = \SI{656,68}{\molar}\)}

    \begin{solu}
%        \begin{figure}[H]
%            \centering
%            \includegraphics[scale=0.8]{images/imagesUD6EquilibrioQuimico/UD6_Problemas/UD6_EVAU_07}
%        \end{figure}

        La reacción es \(\ch{2 NO \gas{} + 2 H2 \gas{} <=> N2 \gas{} + 2 H2O \gas{}}\), en un recipiente de \(V = \SI{1}{\litre}\) (las concentraciones coinciden numéricamente con los moles).

        \textbf{a) Concentraciones en el equilibrio}

        Si llamamos \(2\xi\) a los moles de \ch{NO} que reaccionan, como \([\ch{NO}]_{eq} = 0,062\ \si{\Molar}\):

        \begin{equation*}
            0,1 - 2\xi = 0,062 \ \Rightarrow \ 2\xi = 0,038 \ \Rightarrow \ \xi = 0,019
        \end{equation*}

        \begin{equation*}
            [\ch{H2}]_{eq} = 0,05 - 2\xi = 0,05 - 0,038 = \SI{0.012}{\Molar}
        \end{equation*}

        \begin{equation*}
            [\ch{N2}]_{eq} = \xi = \SI{0.019}{\Molar} \qquad [\ch{H2O}]_{eq} = 0,1 + 2\xi = 0,1 + 0,038 = \SI{0.138}{\Molar}
        \end{equation*}

        \textbf{b) Constante \(K_c\)}

        \begin{equation*}
            K_c = \dfrac{[\ch{N2}][\ch{H2O}]^2}{[\ch{NO}]^2[\ch{H2}]^2} = \dfrac{0,019 \times (0,138)^2}{(0,062)^2 \times (0,012)^2} \approx 654 \ \si{\per\Molar}
        \end{equation*}

    \end{solu}

\end{exerciseEVAU}

\begin{exerciseEVAU}
    La solubilidad del hidróxido de manganeso (II) en agua es de \concentracion[mgl]{1,96}.
    Calcule la constante del producto de solubilidad de dicha sustancia.

    \datos{}: masas atómicas: Mn = 54,94; O = 16; H = 1

% \solution{\(k_{ps} = \num{4,26e-14}\)}

    \begin{solu}
%        \begin{figure}[H]
%            \centering
%            \includegraphics[scale=0.8]{images/imagesUD6EquilibrioQuimico/UD6_Problemas/UD6_EVAU_08}
%        \end{figure}

        Expresamos la solubilidad en mol/L. Con \(M(\ch{Mn(OH)2}) = 54,94 + 2(16+1) = \SI{88.94}{\gram\per\mole}\):

        \begin{equation*}
            s = \dfrac{1,96 \times 10^{-3} \ \si{\gram\per\litre}}{88,94 \ \si{\gram\per\mole}} = \num{2.20e-5} \ \si{\Molar}
        \end{equation*}

        El equilibrio de solubilidad es:

        \begin{equation*}
            \ch{Mn(OH)2 \sld{} <=> Mn^2+ \aq{} + 2 OH- \aq{}}
        \end{equation*}

        Con \([\ch{Mn^2+}] = s\) y \([\ch{OH-}] = 2s\):

        \begin{equation*}
            K_{ps} = [\ch{Mn^2+}][\ch{OH-}]^2 = s\, (2s)^2 = 4 s^3 = 4 \times (2,20 \times 10^{-5})^3 = \num{4.28e-14}
        \end{equation*}

    \end{solu}

\end{exerciseEVAU}

\begin{exerciseEVAU}
    Se coloca cierta cantidad de \ch{SO3} en un matraz de \volumen[l]{0,80}.
    A cierta temperatura se establece el equilibrio de disociación:
    \begin{center}
        \formula{2 SO3\ (g) <=> 2 SO2\ (g) + O2\ (g)}
    \end{center}
    Se comprueba que en el equilibrio hay 2 moles de \ch{O2}.
    Si \(k_c\) es \concentracion[m]{0,22} a la temperatura de la experiencia,
    \begin{enumerate}[label=\letraminuscula]
        \item Calcule las concentraciones de las sustancias presentes en el equilibrio.
        \item Calcule el grado de disociación del \ch{SO3}.
    \end{enumerate}

% \solution{a) \(\conc{SO3} = \concentracion[m]{19,9}\); \(\conc{SO2} = \concentracion[m]{5}\); \(\conc{O2} = \concentracion[m]{2,5}\). b) \(\alpha = 0.23\)}

    \begin{solu}
%        \begin{figure}[H]
%            \centering
%            \includegraphics[scale=0.8]{images/imagesUD6EquilibrioQuimico/UD6_Problemas/UD6_EVAU_09}
%        \end{figure}

        La reacción de disociación es \(\ch{2 SO3 \gas{} <=> 2 SO2 \gas{} + O2 \gas{}}\), en \(V = \SI{0.80}{\litre}\), con \(n(\ch{O2})_{eq} = 2\) mol.

        \textbf{a) Concentraciones en el equilibrio}

        Como se han formado 2 mol de \ch{O2}, y por estequiometría se forman 2 mol de \ch{SO2} por cada mol de \ch{O2}:

        \begin{equation*}
            [\ch{O2}]_{eq} = \dfrac{2}{0,80} = \SI{2.5}{\Molar} \qquad [\ch{SO2}]_{eq} = \dfrac{4}{0,80} = \SI{5}{\Molar}
        \end{equation*}

        Con estos valores y \(K_c = 0,22\ \si{\Molar}\), obtenemos la concentración de \ch{SO3} en el equilibrio:

        \begin{equation*}
            K_c = \dfrac{[\ch{SO2}]^2[\ch{O2}]}{[\ch{SO3}]^2} \ \Rightarrow \ [\ch{SO3}]^2 = \dfrac{[\ch{SO2}]^2[\ch{O2}]}{K_c} = \dfrac{(5)^2 (2,5)}{0,22} = 284,1
        \end{equation*}

        \begin{equation*}
            [\ch{SO3}]_{eq} = \sqrt{284,1} = \SI{16.86}{\Molar}
        \end{equation*}

        \textbf{b) Grado de disociación}

        La cantidad inicial de \ch{SO3} es la que queda en el equilibrio más la que se ha disociado:

        \begin{equation*}
            n_0(\ch{SO3}) = n(\ch{SO3})_{eq} + n(\ch{SO3})_{disociado} = (16,86 \times 0,80) + 4 = 13,49 + 4 = 17,49 \ \text{mol}
        \end{equation*}

        \begin{equation*}
            \alpha = \dfrac{n_{disociado}}{n_0} = \dfrac{4}{17,49} = 0,229 \ (22,9\ \%)
        \end{equation*}

    \end{solu}

\end{exerciseEVAU}

\begin{exerciseEVAU}
    Calcule el producto de solubilidad del \ch{Cd(OH)2} sabiendo que la solubilidad de esta sal en agua es \concentracion[m]{1,4e-5}.

% \solution{\(k_{ps} = \num{1,1e-14}\)}

    \begin{solu}
%        \begin{figure}[H]
%            \centering
%            \includegraphics[scale=0.8]{images/imagesUD6EquilibrioQuimico/UD6_Problemas/UD6_EVAU_10}
%        \end{figure}

        El equilibrio de solubilidad es:

        \begin{equation*}
            \ch{Cd(OH)2 \sld{} <=> Cd^2+ \aq{} + 2 OH- \aq{}}
        \end{equation*}

        Con \(s = \num{1.4e-5}\ \si{\Molar}\), \([\ch{Cd^2+}] = s\) y \([\ch{OH-}] = 2s\):

        \begin{equation*}
            K_{ps} = [\ch{Cd^2+}][\ch{OH-}]^2 = s\, (2s)^2 = 4 s^3 = 4 \times (1,4 \times 10^{-5})^3 = \num{1.10e-14}
        \end{equation*}

    \end{solu}


\end{exerciseEVAU}

\begin{exerciseEVAU}
    En un recipiente de \volumen[l]{2} se introducen \moles{0,40} de \ch{COCl2} y se calienta a \temperatura[k]{900}, con lo que se establece el equilibrio:
    \begin{center}
        \ch{COCl2\ (g) <=> CO\ (g) + Cl2\ (g)}
    \end{center}
    Sabiendo que en ese momento la concentración de \ch{Cl2} es \concentracion[molar]{0,094}:
    \begin{enumerate}[label=\letraminuscula]
        \item Calcule el valor del grado de disociación del \ch{COCl2}.
        \item Calcule el valor de \(kc\)  y \(kp\) .
        \item Explique cómo afectaría a la concentración de \ch{COCl2} en la mezcla gaseosa en equilibrio la adición de \moles{0,2} de \ch{Cl2} manteniendo constante la temperatura.
        (No se requiere cálculo numérico)
    \end{enumerate}
    \datos{} R = \ctegases{}

% \solution{a) \(\alpha = 0.47\). b) \(k_c = \SI{0,083}{\molar}\). c) Aumenta la concentración de \ch{COCl2}}

    \begin{solu}
%        \begin{figure}[H]
%            \centering
%            \includegraphics[scale=0.8]{images/imagesUD6EquilibrioQuimico/UD6_Problemas/UD6_EVAU_11}
%        \end{figure}

        La reacción es \(\ch{COCl2 \gas{} <=> CO \gas{} + Cl2 \gas{}}\), con \(V = \SI{2}{\litre}\) y \(T = \SI{900}{\kelvin}\). La concentración inicial es \([\ch{COCl2}]_0 = 0,40/2 = \SI{0.20}{\Molar}\).

        \textbf{a) Grado de disociación}

        Como \([\ch{Cl2}]_{eq} = x = \SI{0.094}{\Molar}\):

        \begin{equation*}
            \alpha = \dfrac{x}{[\ch{COCl2}]_0} = \dfrac{0,094}{0,20} = 0,47 \ (47\ \%)
        \end{equation*}

        \textbf{b) \(K_c\) y \(K_p\)}

        \begin{equation*}
            [\ch{COCl2}]_{eq} = 0,20 - 0,094 = \SI{0.106}{\Molar}
        \end{equation*}

        \begin{equation*}
            K_c = \dfrac{[\ch{CO}][\ch{Cl2}]}{[\ch{COCl2}]} = \dfrac{(0,094)^2}{0,106} = \SI{0.083}{\Molar}
        \end{equation*}

        Como \(\Delta n = 2 - 1 = 1\):

        \begin{equation*}
            K_p = K_c\, (RT) = 0,083 \times (0,082 \times 900) = \SI{6.15}{\atmospher}
        \end{equation*}

        \textbf{c) Efecto de añadir 0,2 mol de \ch{Cl2}}

        Al añadir \ch{Cl2}, aumenta la concentración de uno de los productos, por lo que, según el principio de Le Châtelier, el equilibrio se desplaza hacia la izquierda (hacia los reactivos), consumiendo parte del \ch{Cl2} añadido y del \ch{CO} presente. Esto provoca que \textbf{aumente la concentración de \ch{COCl2}} en la mezcla en equilibrio (sin necesidad de cálculo numérico, ya que \(K_c\) se mantiene constante a temperatura constante).

    \end{solu}

\end{exerciseEVAU}

\begin{exerciseEVAU}

    La solubilidad del \ch{Ag2CrO4} en agua a \temperatura[c]{25} es \concentracion[molar]{1e-4}, calcule el producto de solubilidad de esta sal.

% \solution{\(k_{ps} = \num{4e-12}\)}

    \begin{solu}
%        \begin{figure}[H]
%            \centering
%            \includegraphics[scale=0.8]{images/imagesUD6EquilibrioQuimico/UD6_Problemas/UD6_EVAU_12}
%        \end{figure}

        El equilibrio de solubilidad es:

        \begin{equation*}
            \ch{Ag2CrO4 \sld{} <=> 2 Ag+ \aq{} + CrO4^2- \aq{}}
        \end{equation*}

        Con \(s = \num{1e-4}\ \si{\Molar}\), \([\ch{Ag+}] = 2s\) y \([\ch{CrO4^2-}] = s\):

        \begin{equation*}
            K_{ps} = [\ch{Ag+}]^2[\ch{CrO4^2-}] = (2s)^2\, s = 4 s^3 = 4 \times (1 \times 10^{-4})^3 = \num{4e-12}
        \end{equation*}

    \end{solu}

\end{exerciseEVAU}

\begin{exerciseEVAU}
    El alcohol etílico y el ácido acético reaccionan de acuerdo con la ecuación:
    \begin{center}
        \ch{C2H5OH + CH3COOH <=> CH3COOC2H5 + H2O}
    \end{center}
    Si en un recipiente de 1 litro se mezcla 1 mol de etanol y 1 mol de ácido acético se alcanza el equilibrio cuando se han formado 2/3 de moles de \ch{CH3COOC2H5} y 2/3 de moles de agua.
    \begin{enumerate}[label=\letraminuscula]
\item Calcule la constante de equilibrio para la anterior ecuación de reacción.
\item La presión total cuando se alcanza el equilibrio a \temperatura[c]{200}.
\item La composición del equilibrio al mezclar \moles{1} de \ch{CH3COOC2H5} y \moles{1} de agua
    \end{enumerate}

    \begin{solu}
        \begin{table}[H]
            % --- tabla antigua ---
            % \begin{tblr}{%
                % colspec = {X[c]X[c]X[c]X[c]X[c]},
                % vline{2} = {1}{ text = \clap{:} },
                % vline{3} = {1}{ text = \clap{\ch{+}} },
                % vline{4} = {1}{ text = \clap{\ch{->}} },
                % vline{5} = {1}{ text = \clap{\ch{+}} },
            % }
                % Ecuación         & \ch{C2H5OH}        & \ch{CH3COOH}       & \ch{CH3COOC2H5} & \ch{H2O}         \\
                % \(n_0\)          & \(1\)              & \(1\)              & \(-\)           & \(-\)            \\
                % \(n\) reaccionan & \(-x\)             & \(-x\)             & \(-\)           & \(-\)            \\
                % \(n\) forman     & \(-\)              & \( - \)            & \(+x \)         & \(+x\)           \\
                % \(n\) equilibrio & \(1-x\)            & \(1-x\)            & \(x\)           & \(x\)            \\
                % \(eq\)           & \(\dfrac{1-x}{1}\) & \(\dfrac{1-x}{1}\) & \(\dfrac{x}{1}\) & \(\dfrac{x}{1}\) \\
            % \end{tblr}
            % --- tabla nueva ---
            \begin{tabular}{@{}c @{\,}c@{\,} c @{\,}c@{\,} c @{\,}c@{\,} c @{\,}c@{\,} c@{}}
                \toprule
                \textbf{Ecuación} & : & \textbf{\ch{C2H5OH}} & \ch{+} & \textbf{\ch{CH3COOH}} & \ch{->} & \textbf{\ch{CH3COOC2H5}} & \ch{+} & \textbf{\ch{H2O}} \\
                \midrule
                \(n_0\) &  & \(1\) &  & \(1\) &  & \(-\) &  & \(-\) \\[3pt]
                \(n\) reaccionan &  & \(-x\) &  & \(-x\) &  & \(-\) &  & \(-\) \\[3pt]
                \(n\) forman &  & \(-\) &  & \( - \) &  & \(+x \) &  & \(+x\) \\[3pt]
                \(n\) equilibrio &  & \(1-x\) &  & \(1-x\) &  & \(x\) &  & \(x\) \\[3pt]
                \(eq\) &  & \(\dfrac{1-x}{1}\) &  & \(\dfrac{1-x}{1}\) &  & \(\dfrac{x}{1}\) &  & \(\dfrac{x}{1}\) \\
                \bottomrule
            \end{tabular}
        \end{table}

        Como \(V = \SI{1}{\litre}\), las concentraciones coinciden numéricamente con los moles.

        \begin{enumerate}[label=\letraminuscula]
\item Sabemos que en el equilibrio se han formado \(x = 2/3\) mol de \ch{CH3COOC2H5} y de \ch{H2O}, por lo que:
\begin{equation*}
    n(\ch{C2H5OH})_{eq} = n(\ch{CH3COOH})_{eq} = 1 - \dfrac{2}{3} = \dfrac{1}{3} \ \text{mol}
\end{equation*}
Las concentraciones en el equilibrio son, por tanto, \([\ch{C2H5OH}] = [\ch{CH3COOH}] = \SI{0.333}{\Molar}\) y \([\ch{CH3COOC2H5}] = [\ch{H2O}] = \SI{0.667}{\Molar}\). La constante de equilibrio vale:
\begin{equation*}
    k_c = \dfrac{\conc{CH3COOC2H5}\ \conc{H2O}}{\conc{C2H5OH}\ \conc{CH3COOH}} = \dfrac{(2/3)(2/3)}{(1/3)(1/3)} = \dfrac{4/9}{1/9} = 4
\end{equation*}
\item Como el número de moles gaseosos no cambia en la reacción (2 mol de reactivos dan 2 mol de productos, \(\Delta n = 0\)), los moles totales en el equilibrio son siempre \(n_T = 1+1 = 2\) mol, independientemente del grado de avance. A \(T = \SI{473}{\kelvin}\) (\(200+273\)) y \(V = \SI{1}{\litre}\):
\begin{equation*}
    p = \dfrac{n_T\, R\, T}{V} = \dfrac{2 \times 0,082 \times 473}{1} = \SI{77.57}{\atmospher}
\end{equation*}
\item Partimos ahora de 1 mol de \ch{CH3COOC2H5} y 1 mol de \ch{H2O} (los productos), y llamamos \(y\) a los moles que reaccionan en sentido inverso, formando \ch{C2H5OH} y \ch{CH3COOH}:
\begin{equation*}
    n(\ch{CH3COOC2H5})_{eq} = n(\ch{H2O})_{eq} = 1-y \qquad n(\ch{C2H5OH})_{eq} = n(\ch{CH3COOH})_{eq} = y
\end{equation*}
Como la constante de equilibrio es la misma (\(k_c = 4\)), sustituyendo en su expresión:
\begin{equation*}
    k_c = \dfrac{(1-y)^2}{y^2} = 4 \ \Rightarrow \ \dfrac{1-y}{y} = 2 \ \Rightarrow \ 1-y = 2y \ \Rightarrow \ y = \dfrac{1}{3}
\end{equation*}
Por tanto, la composición en el equilibrio es:
\begin{equation*}
    [\ch{CH3COOC2H5}] = [\ch{H2O}] = 1 - \dfrac{1}{3} = \SI{0.667}{\Molar} \qquad [\ch{C2H5OH}] = [\ch{CH3COOH}] = \dfrac{1}{3} = \SI{0.333}{\Molar}
\end{equation*}
Como cabía esperar, al ser la reacción de estequiometría simétrica (1:1:1:1) con la misma cantidad total de moles (2) en ambos casos, se alcanza exactamente la misma composición de equilibrio partiendo de reactivos puros o de productos puros.
        \end{enumerate}

    \end{solu}

\end{exerciseEVAU}

\begin{exerciseEVAU}
    Se mezclan \volumen[ml]{10} de \ch{BaCl2} \concentracion[m]{0,1} con \volumen[ml]{40} de \ch{Na2SO4} \concentracion[m]{0,1}.

    \datos{} \(K_{ps} \ch{BaSO4} = \num{1,1e-10}\)
    \begin{enumerate}[label=\letraminuscula]
\item ¿Precipitará \ch{BaSO4}?
\item Escriba el equilibrio de solubilidad del \ch{BaSO4}
    \end{enumerate}

    \begin{solu}
%        \begin{figure}[H]
%            \centering
%            \includegraphics[scale=0.8]{images/imagesUD6EquilibrioQuimico/UD6_Problemas/UD6_EVAU_14}
%        \end{figure}

        \textbf{a) ¿Precipitará \ch{BaSO4}?}

        Calculamos los moles de cada ion aportados y, tras la mezcla, sus concentraciones en el volumen total:

        \begin{equation*}
            n(\ch{Ba^2+}) = 0,010 \ \si{\litre} \times 0,1 \ \si{\Molar} = \num{1e-3} \ \text{mol} \qquad n(\ch{SO4^2-}) = 0,040 \ \si{\litre} \times 0,1 \ \si{\Molar} = \num{4e-3} \ \text{mol}
        \end{equation*}

        El volumen total es \(V = 10 + 40 = \SI{50}{\milli\litre} = \SI{0.050}{\litre}\), por lo que:

        \begin{equation*}
            [\ch{Ba^2+}] = \dfrac{10^{-3}}{0,050} = \SI{0.02}{\Molar} \qquad [\ch{SO4^2-}] = \dfrac{4 \times 10^{-3}}{0,050} = \SI{0.08}{\Molar}
        \end{equation*}

        \begin{equation*}
            Q = [\ch{Ba^2+}][\ch{SO4^2-}] = 0,02 \times 0,08 = \num{1.6e-3}
        \end{equation*}

        Como \(Q = 1,6 \times 10^{-3} \gg K_{ps} = 1,1 \times 10^{-10}\), la disolución está sobresaturada y \textbf{sí precipitará} \ch{BaSO4}.

        \textbf{b) Equilibrio de solubilidad}

        \begin{equation*}
            \ch{BaSO4 \sld{} <=> Ba^2+ \aq{} + SO4^2- \aq{}} \qquad K_{ps} = [\ch{Ba^2+}][\ch{SO4^2-}]
        \end{equation*}

    \end{solu}

% \solution{a) Si precipitará el \ch{BaSO4}. b) \(k_{ps} = \conc{Ba^{2+}}\ \conc{SO4^{2-}}\)}

\end{exerciseEVAU}

\begin{exerciseEVAU}
    En un recipiente de \volumen[l]{2,0} de capacidad se introduce una mezcla gaseosa que contiene \moles{0,1} de \ch{SO2} y \num{0,5} de \ch{SO2Cl2} (cloruro de sulfonilo). A \temperatura[c]{150} se establece el equilibrio :
    \begin{center}
        \ch{SO2Cl2\ (g) <=> SO2\ (g) + Cl2\ (g)}
    \end{center}
    cuya constante \(Kc\) vale \num{0,011}.
    \begin{enumerate}[label=\letraminuscula]
\item Calcular la concentración de todas las sustancias en equilibrio.
\item Calcular la presión total en el equilibrio a \temperatura[c]{150}
\item Calcular Kp y los valores de las presiones parciales de los componentes de la mezcla en equilibrio.
    \end{enumerate}

    \datos{}: R = \ctegases

% \solution{a) \(\conc{SO2Cl2} = \concentracion[m]{0,22}\); \(\conc{SO2} = \concentracion[m]{0,08}\); \(\conc{Cl2} = \concentracion[m]{0,03}\). b) \(p=\SI{11,45}{\atmospher}\). c) \(k_p = \SI{0,38}{\atmospher}\); \(p_{\ch{SO2Cl2}} = \SI{7,63}{\atmospher}\); \(p_{\ch{SO2}} = \SI{2,78}{\atmospher}\); \(p_{\ch{Cl2}} = \SI{1,04}{\atmospher}\) }

    \begin{solu}
%        \begin{figure}[H]
%            \centering
%            \includegraphics[scale=0.8]{images/imagesUD6EquilibrioQuimico/UD6_Problemas/UD6_EVAU_15A}
%            \includegraphics[scale=0.8]{images/imagesUD6EquilibrioQuimico/UD6_Problemas/UD6_EVAU_15B}
%        \end{figure}

        La reacción es \(\ch{SO2Cl2 \gas{} <=> SO2 \gas{} + Cl2 \gas{}}\), con \(V=\SI{2.0}{\litre}\) y \(T = \SI{423}{\kelvin}\) (\(150+273\)). Las concentraciones iniciales son \([\ch{SO2Cl2}]_0 = 0,5/2 = \SI{0.25}{\Molar}\) y \([\ch{SO2}]_0 = 0,1/2 = \SI{0.05}{\Molar}\).

        \textbf{a) Concentraciones en el equilibrio}

        Llamando \(x\) a la concentración de \ch{SO2Cl2} que se disocia:

        \begin{equation*}
            K_c = \dfrac{(0,05+x)\, x}{0,25 - x} = 0,011
        \end{equation*}

        Desarrollando: \(x^2 + 0,061\, x - 0,00275 = 0\), cuya raíz positiva es \(x = 0,0302\ \si{\Molar}\). Por tanto:

        \begin{equation*}
            [\ch{SO2Cl2}]_{eq} = 0,25 - 0,0302 \approx \SI{0.22}{\Molar} \qquad [\ch{SO2}]_{eq} = 0,05 + 0,0302 \approx \SI{0.08}{\Molar} \qquad [\ch{Cl2}]_{eq} \approx \SI{0.03}{\Molar}
        \end{equation*}

        \textbf{b) Presión total en el equilibrio}

        \begin{equation*}
            n_T = ([\ch{SO2Cl2}]+[\ch{SO2}]+[\ch{Cl2}]) \times V = (0,22+0,08+0,03) \times 2 = 0,66 \ \text{mol}
        \end{equation*}

        \begin{equation*}
            P_T = \dfrac{n_T\, R\, T}{V} = \dfrac{0,66 \times 0,082 \times 423}{2} = \SI{11.45}{\atmospher}
        \end{equation*}

        \textbf{c) \(K_p\) y presiones parciales}

        Como \(\Delta n = 2-1 = 1\):

        \begin{equation*}
            K_p = K_c\, (RT) = 0,011 \times (0,082 \times 423) = \SI{0.38}{\atmospher}
        \end{equation*}

        Las presiones parciales se obtienen con \(p_i = [i]\, R\, T\):

        \begin{equation*}
            p_{\ch{SO2Cl2}} = 0,22 \times 34,69 = \SI{7.63}{\atmospher} \qquad p_{\ch{SO2}} = 0,08 \times 34,69 = \SI{2.78}{\atmospher} \qquad p_{\ch{Cl2}} = 0,03 \times 34,69 = \SI{1.04}{\atmospher}
        \end{equation*}

    \end{solu}

\end{exerciseEVAU}

\begin{exerciseEVAU}
    El \ch{Al(OH)3} es una sustancia muy poco soluble en agua (\(Ks=\num{3,7e-15}\)).
    \begin{enumerate}[label=\letraminuscula]
\item Escribir el equilibrio de solubilidad y la expresión del producto de solubilidad.
\item ¿Cuál es su solubilidad en agua pura, expresada en \si{\g\per\litre}?
\item ¿Cómo afectaría a la solubilidad del \ch{Al(OH)3} una disminución del \(pH\) de la disolución? Razone en términos del principio de Le Chatelier-Braun.
    \end{enumerate}

    \datos{}: Masas atómicas Al = 27,0; O = 16,0.

% \solution{\(k_{ps} = 27\ s^4\). b) \(s = \SI{8,6e-3}{\g\per\litre}\). c) aumentará la solubilidad}

    \begin{solu}
%        \begin{figure}[H]
%            \centering
%            \includegraphics[scale=0.8]{images/imagesUD6EquilibrioQuimico/UD6_Problemas/UD6_EVAU_16A}
%        \end{figure}

        \textbf{a) Equilibrio de solubilidad}

        \begin{equation*}
            \ch{Al(OH)3 \sld{} <=> Al^3+ \aq{} + 3 OH- \aq{}} \qquad K_s = [\ch{Al^3+}][\ch{OH-}]^3
        \end{equation*}

        \textbf{b) Solubilidad en agua pura}

        Si \(s\) es la solubilidad, \([\ch{Al^3+}] = s\) y \([\ch{OH-}] = 3s\):

        \begin{equation*}
            K_s = s\, (3s)^3 = 27 s^4 = \num{3.7e-15} \ \Rightarrow \ s^4 = \num{1.37e-16} \ \Rightarrow \ s = \num{1.08e-4} \ \si{\Molar}
        \end{equation*}

        Con \(M(\ch{Al(OH)3}) = 27,0 + 3(16,0+1,0) = \SI{78.0}{\gram\per\mole}\):

        \begin{equation*}
            s = 1,08 \times 10^{-4} \ \text{mol/L} \times 78,0 \ \si{\gram\per\mole} = \num{8.44e-3} \ \si{\gram\per\litre}
        \end{equation*}

        \textbf{c) Efecto de una disminución del pH}

        Al disminuir el pH, aumenta \([\ch{H+}]\), que reacciona con el \ch{OH-} presente en disolución (neutralización), disminuyendo \([\ch{OH-}]\). Según el principio de Le Châtelier-Braun, el equilibrio de solubilidad se desplaza hacia la derecha (más disolución del sólido) para reponer parte del \ch{OH-} consumido, por lo que \textbf{aumenta la solubilidad} del \ch{Al(OH)3}.

    \end{solu}

\end{exerciseEVAU}

\begin{exerciseEVAU}
    Para la reacción química \ch{A + B -> 2 C}, se conoce el valor absoluto de la variación de entalpía, \(|\Delta H| = \energia[kj]{68}\) y la variación de entropía, \(\Delta S = \SI{2}{\joule\per\kelvin}\).
    \begin{enumerate}[label=\letraminuscula]
\item Comente razonadamente si la reacción es espontánea a \temperatura[c]{25}.
\item Si la ecuación de velocidad de reacción tiene la forma \(v = k\ \conc{A}^2\ \conc{B}\), ¿cómo afecta a la velocidad reducir a la mitad la concentración de A al tiempo que se duplica la de B?
    \end{enumerate}

    \begin{solu}
%        \begin{figure}[H]
%            \centering
%            \includegraphics[scale=0.8]{images/imagesUD6EquilibrioQuimico/UD6_Problemas/UD6_EVAU_17}
%        \end{figure}

        \textbf{a) Espontaneidad a 25 ºC}

        A \(T = \SI{298}{\kelvin}\), el término entrópico vale:

        \begin{equation*}
            T \Delta S = 298 \times 2 = \SI{596}{\joule} = \SI{0.596}{\kilo\joule}
        \end{equation*}

        Este valor es muy pequeño frente a \(|\Delta H| = \SI{68}{\kilo\joule}\), por lo que la espontaneidad depende, prácticamente en su totalidad, del signo de \(\Delta H\):

        \begin{itemize}
            \item Si la reacción es \textbf{exotérmica} (\(\Delta H = \energia[kj]{-68}\)): \(\Delta G = -68 - 0,596 = \energia[kj]{-68.6} < 0\), \textbf{espontánea}.
            \item Si la reacción es \textbf{endotérmica} (\(\Delta H = \energia[kj]{68}\)): \(\Delta G = 68 - 0,596 = \energia[kj]{67.4} > 0\), \textbf{no espontánea}.
        \end{itemize}

        \textbf{b) Efecto sobre la velocidad de reacción}

        Con \(v = k\, [A]^2\, [B]\), si \([A]\) se reduce a la mitad y \([B]\) se duplica:

        \begin{equation*}
            v' = k\, \left( \dfrac{[A]}{2} \right)^2 (2[B]) = k\, \dfrac{[A]^2}{4}\, (2[B]) = \dfrac{2}{4}\, k\, [A]^2 [B] = \dfrac{1}{2}\, v
        \end{equation*}

        La velocidad de reacción se \textbf{reduce a la mitad}.

    \end{solu}

\end{exerciseEVAU}

\begin{exerciseEVAU}
    En un recipiente de 2 litros se introduce inicialmente una mezcla de 1,5 moles del compuesto gaseoso A y 2,5 moles del compuesto gaseoso B y se calienta a \temperatura[k]{500}, estableciéndose el equilibrio indicado por la reacción:
    \begin{center}
        \ch{A\ (g) + 2 B\ (g) <=> 2 C\ (g)}
    \end{center}
    Sabiendo que la fracción molar del compuesto C en la mezcla en equilibrio es 0,35 calcula:
    \begin{enumerate}[label=\letraminuscula]
\item Las concentraciones molares de todos los compuestos en el equilibrio.
\item El valor de las constantes de equilibrio \(K_c\) y \(K_p\).
\item La presión total en el recipiente cuando se alcanza el equilibrio a \temperatura[k]{500}.
    \end{enumerate}
    \datos{}: R = \ctegases

    \begin{solu}
%        \begin{figure}[H]
%            \centering
%            \includegraphics[scale=0.8]{images/imagesUD6EquilibrioQuimico/UD6_Problemas/UD6_EVAU_18}
%        \end{figure}

        La reacción es \(\ch{A \gas{} + 2 B \gas{} <=> 2 C \gas{}}\), con \(V = \SI{2}{\litre}\) y \(T = \SI{500}{\kelvin}\). Llamando \(\xi\) a los moles de A que reaccionan:

        \begin{equation*}
            n(A) = 1,5-\xi \qquad n(B) = 2,5-2\xi \qquad n(C) = 2\xi \qquad n_T = 4 - \xi
        \end{equation*}

        \textbf{a) Concentraciones en el equilibrio}

        Con \(x_C = \dfrac{2\xi}{4-\xi} = 0,35\):

        \begin{equation*}
            2\xi = 0,35(4-\xi) \ \Rightarrow \ 2,35\, \xi = 1,4 \ \Rightarrow \ \xi = 0,596 \ \text{mol}
        \end{equation*}

        \begin{equation*}
            n(A) = 0,904 \ \text{mol} \qquad n(B) = 1,309 \ \text{mol} \qquad n(C) = 1,191 \ \text{mol} \qquad n_T = 3,404 \ \text{mol}
        \end{equation*}

        \begin{equation*}
            [A] = \dfrac{0,904}{2} = \SI{0.452}{\Molar} \qquad [B] = \dfrac{1,309}{2} = \SI{0.654}{\Molar} \qquad [C] = \dfrac{1,191}{2} = \SI{0.596}{\Molar}
        \end{equation*}

        \textbf{b) Constantes \(K_c\) y \(K_p\)}

        \begin{equation*}
            K_c = \dfrac{[C]^2}{[A][B]^2} = \dfrac{(0,596)^2}{0,452 \times (0,654)^2} = 1,83
        \end{equation*}

        Como \(\Delta n = 2-3=-1\):

        \begin{equation*}
            K_p = K_c\, (RT)^{-1} = \dfrac{1,83}{0,082 \times 500} = 0,0447 \ \si{\per\atmospher}
        \end{equation*}

        \textbf{c) Presión total}

        \begin{equation*}
            P_T = \dfrac{n_T\, R\, T}{V} = \dfrac{3,404 \times 0,082 \times 500}{2} = \SI{69.8}{\atmospher}
        \end{equation*}

    \end{solu}


\end{exerciseEVAU}

\begin{exerciseEVAU}
    El diyoduro de plomo es un compuesto muy poco soluble en agua.
    \begin{enumerate}[label=\letraminuscula]
\item Escribe el equilibrio de solubilidad del diyoduro de plomo.
\item Dónde será más soluble el diyoduro de plomo: en agua pura o en una disolución 0,1 M de yoduro de potasio? Justifica la respuesta en términos del principio de Le Chatelier
    \end{enumerate}

    \begin{solu}
%        \begin{figure}[H]
%            \centering
%            \includegraphics[scale=0.8]{images/imagesUD6EquilibrioQuimico/UD6_Problemas/UD6_EVAU_19}
%        \end{figure}

        \textbf{a) Equilibrio de solubilidad}

        \begin{equation*}
            \ch{PbI2 \sld{} <=> Pb^2+ \aq{} + 2 I- \aq{}} \qquad K_s = [\ch{Pb^2+}][\ch{I-}]^2
        \end{equation*}

        \textbf{b) Comparación de solubilidades}

        En la disolución de \ch{KI} 0,1 M ya existe una concentración inicial de ion \ch{I-} (ion común con el equilibrio de solubilidad del \ch{PbI2}). Según el principio de Le Châtelier, la presencia de este exceso de \ch{I-} desplaza el equilibrio de solubilidad hacia la \textbf{izquierda} (hacia el sólido), dificultando la disolución del \ch{PbI2}.

        Por tanto, el \ch{PbI2} será \textbf{más soluble en agua pura} que en la disolución 0,1 M de \ch{KI}, donde su solubilidad se ve reducida por el efecto del ion común.

    \end{solu}

\end{exerciseEVAU}

\begin{exerciseEVAU}
    Por deshidrogenación del metilciclohexano (\ch{C7H14}) a \temperatura[k]{700} se obtiene tolueno (\ch{C7H8}), de acuerdo con el equilibrio:
    \begin{equation*}
        \ch{C7H14\ (g) <=> C7H8\ (g) + 3 H2\ (g)}\; \Delta H > 0
    \end{equation*}

    En un recipiente de 1 litro, inicialmente vacío, se introducen \moles{0,6} de metilciclohexano y se calientan a \temperatura[k]{700}, de forma que, establecido el equilibrio, hay \moles{0,45} de \ch{H2} en la mezcla gaseosa. Calcula:
    \begin{enumerate}[label=\letraminuscula]
\item La constante Kp a dicha temperatura.
\item El grado de disociación del metilciclohexano.
\item ¿Qué efecto tendrá sobre la fracción molar del tolueno en la mezcla un aumento de la temperatura? ¿Y la adición de un catalizador adecuado? Razona las respuestas.
    \end{enumerate}
    \datos{} R = \ctegases

    \begin{solu}

        Primero ponemos el equilibrio:
        \begin{table}[H]
            % --- tabla antigua ---
            % \begin{tblr}{%
                % colspec = {X[c]X[c]X[c]X[c]},
% % cell{3}{2} = {red9},
                % vline{2} = {1}{ text = \clap{:} },
% % vline{3} = {1}{ text = \clap{\ch{+}} },
                % vline{3} = {1}{ text = \clap{\ch{->}} },
                % vline{4} = {1}{ text = \clap{\ch{+}} },
            % }
                % Ecuación             & \ch{C7H14\ (g)}        & \ch{C7H8}\ (g)    & \ch{3 H2}\ (g)      \\
                % \( n_0 \)            & \(0.6\)                & \(-\)             & \(-\)               \\
                % \( n_{reaccionan} \) & \(-x\)                 & \(-\)             & \(-\)               \\
                % \( n_{forman} \)     & \(-\)                  & \(+x\)            & \(+3x\)             \\
                % \( n_{equilibrio}\)  & \( 0.6-x \)            & \( x \)           & \( 3x \)            \\
                % \conc{eq}            & \( \dfrac{0.6-x}{1} \) & \(\dfrac{x}{1} \) & \( \dfrac{3x}{1} \) \\
            % \end{tblr}
            % --- tabla nueva ---
            \begin{tabular}{@{}c @{\,}c@{\,} c @{\,}c@{\,} c @{\,}c@{\,} c@{}}
                \toprule
                \textbf{Ecuación} & : & \textbf{\ch{C7H14\ (g)}} & \ch{->} & \textbf{\ch{C7H8}\ (g)} & \ch{+} & \textbf{\ch{3 H2}\ (g)} \\
                \midrule
                \( n_0 \) &  & \(0.6\) &  & \(-\) &  & \(-\) \\[3pt]
                \( n_{reaccionan} \) &  & \(-x\) &  & \(-\) &  & \(-\) \\[3pt]
                \( n_{forman} \) &  & \(-\) &  & \(+x\) &  & \(+3x\) \\[3pt]
                \( n_{equilibrio}\) &  & \( 0.6-x \) &  & \( x \) &  & \( 3x \) \\[3pt]
                \conc{eq} &  & \( \dfrac{0.6-x}{1} \) &  & \(\dfrac{x}{1} \) &  & \( \dfrac{3x}{1} \) \\
                \bottomrule
            \end{tabular}
        \end{table}

        \begin{enumerate}[label=\letraminuscula]
\item Cuando se alcanza el equilibrio sabemos que tenemos \moles{0.45} de \ch{H2} y por lo tanto tendremos \(3x=0.45\) es decir que \(x=\dfrac{0.45}{3} = \moles{0.15}\).

Ahora, podemos calcular las concentraciones en equilibrio:

\begin{multicols}{2}
    \begin{itemize}
        \item \([\ch{C7H14}] = \dfrac{0.6-0.15}{1} = \SI{0.45}{\molar}\).
        \item \([\ch{C7H8}] = \dfrac{0.15}{1} = \SI{0.15}{\molar}\).
        \item \([\ch{H2}] = \dfrac{0.45}{1} = \SI{0.45}{\molar}\).
    \end{itemize}
\end{multicols}
Primero, calculamos la \(k_c\) según:
\begin{gather*}
    k_c = \dfrac{\conc{H2}^3\ \conc{C7H8}}{\conc{C7H14}}\separ
    k_c = \dfrac{0.45^3 \cdot 0.15}{0.45} = 0.03
\end{gather*}
Para calcular la \(k_p\):
\begin{gather*}
    k_p = k_c\ (R\ T)^{\Delta n} \rightarrow k_p = 0.03\cdot (0.082\cdot 700)^{4-1} \rightarrow k_p = \SI{5673}{\atmospher}
\end{gather*}

Otra posible manera de hacerlo sería, calcular en primer lugar el número total de moles en el equilibrio, una vez obtenidos podríamos sacar la presión total en el equilibrio y a partir de ahí, obtener las presiones parciales de cada componente utilizando la ley de Dalton.

\item Para calcular el grado de disociación:
\begin{gather*}
    \alpha = \dfrac{n_{\text{disociado}}}{n_0} = \dfrac{0.15}{0.6} = 0.25.
\end{gather*}
\item Un aumento de la temperatura, favorecería una aumento en la concentración de tolueno, ya que la reacción se desplazaría hacia donde absorberá energía, es decir hacia la derecha.
Si añadimos un catalizador, las concentraciones se mantendrán inalteradas, pero aumentará la velocidad de reacción.
        \end{enumerate}

    \end{solu}


\end{exerciseEVAU}

\begin{exerciseEVAU}
    El \ch{Fe(OH)3} es un compuesto poco soluble en agua.
    a) Escribe el equilibrio de solubilidad del \ch{Fe(OH)3}.
    b) ¿Cómo afecta a la solubilidad de este compuesto un aumento del \(pH\)? Justifica la respuesta.

    \begin{solu}
%        \begin{figure}[H]
%            \centering
%            \includegraphics[scale=0.8]{images/imagesUD6EquilibrioQuimico/UD6_Problemas/UD6_EVAU_21}
%        \end{figure}

        \textbf{a) Equilibrio de solubilidad}

        \begin{equation*}
            \ch{Fe(OH)3 \sld{} <=> Fe^3+ \aq{} + 3 OH- \aq{}} \qquad K_s = [\ch{Fe^3+}][\ch{OH-}]^3
        \end{equation*}

        \textbf{b) Efecto de un aumento del pH}

        Un aumento del pH implica un aumento de \([\ch{OH-}]\) en la disolución. Según el principio de Le Châtelier, este aumento de la concentración de uno de los iones del equilibrio de solubilidad desplaza dicho equilibrio hacia la \textbf{izquierda} (hacia el sólido, es decir, favoreciendo la precipitación), por lo que \textbf{disminuye la solubilidad} del \ch{Fe(OH)3}.

    \end{solu}

\end{exerciseEVAU}

\begin{exerciseEVAU}
    El\ch{Cd(OH)2} es una sustancia cuyo producto de solubilidad en agua a \temperatura[c]{25} es \num{7,2e-15}\begin{enumerate}[label=\letraminuscula]
\item Escribe el equilibrio de solubilidad del \ch{Cd(OH)2}.
\item Calcula su solubilidad en \si{\g\per\litre} a \temperatura[c]{25}.
\item Razona cómo afectará a la solubilidad de esta sustancia una reducción del pH del medio.
    \end{enumerate}
    \datos{}Masas atómicas: H = 1,0; = = 16,0; Cd = 112,4

    \begin{solu}
%        \begin{figure}[H]
%            \centering
%            \includegraphics[scale=0.8]{images/imagesUD6EquilibrioQuimico/UD6_Problemas/UD6_EVAU_22}
%        \end{figure}

        \textbf{a) Equilibrio de solubilidad}

        \begin{equation*}
            \ch{Cd(OH)2 \sld{} <=> Cd^2+ \aq{} + 2 OH- \aq{}} \qquad K_s = [\ch{Cd^2+}][\ch{OH-}]^2
        \end{equation*}

        \textbf{b) Solubilidad en g/L}

        Si \(s\) es la solubilidad molar, \([\ch{Cd^2+}] = s\) y \([\ch{OH-}] = 2s\):

        \begin{equation*}
            K_s = s\, (2s)^2 = 4 s^3 = \num{7.2e-15} \ \Rightarrow \ s^3 = \num{1.8e-15} \ \Rightarrow \ s = \num{1.22e-5} \ \si{\Molar}
        \end{equation*}

        Con \(M(\ch{Cd(OH)2}) = 112,4 + 2(16,0+1,0) = \SI{146.4}{\gram\per\mole}\):

        \begin{equation*}
            s = 1,22 \times 10^{-5} \ \text{mol/L} \times 146,4 \ \si{\gram\per\mole} = \num{1.78e-3} \ \si{\gram\per\litre}
        \end{equation*}

        \textbf{c) Efecto de una reducción del pH}

        Al reducir el pH aumenta \([\ch{H+}]\), que neutraliza al \ch{OH-} presente en disolución. Según el principio de Le Châtelier, el equilibrio de solubilidad se desplaza hacia la derecha (más disolución) para reponer parte del \ch{OH-} consumido, por lo que \textbf{aumenta la solubilidad} del \ch{Cd(OH)2}.

    \end{solu}


\end{exerciseEVAU}

\begin{exerciseEVAU}
    En un recipiente de \volumen[l]{5,00} se introducen un mol de dióxido de azufre y un mol de oxígeno, y se calienta el sistema a \temperatura[c]{1000} con lo que se produce la reacción: \ch{2 SO2\ (g) + O2\ (g) <=> 2 SO3\ (g)}. Sabiendo que en el equilibrio hay \moles{0,15} de dióxido de azufre, calcula:
    \begin{enumerate}[label=\letraminuscula]
\item Las concentraciones molares de todos los compuestos en el equilibrio.
\item Los valores de \(K_c\) y \(K_p\) .
\item ¿Cómo afecta a la concentración de \ch{SO3} en la mezcla gaseosa un aumento de la presión en el recipiente manteniendo constante la temperatura? Razona la respuesta.
    \end{enumerate}
    \datos{} R = \ctegases

    \begin{solu}
        \begin{table}[H]
            % --- tabla antigua ---
            % \begin{tblr}{%
                % colspec = {X[c]X[c]X[c]X[c]},
% % cell{3}{2} = {red9},
                % vline{2} = {1}{ text = \clap{:} },
                % vline{3} = {1}{ text = \clap{\ch{+}} },
                % vline{4} = {1}{ text = \clap{\ch{->}} },
            % }
                % \textbf{Ecuación}              & \ch{2 SO2\ (g)}                                                         & \ch{O2 \ (g)}                                                         & \ch{2 SO3\ (g)}                                                     \\
                % $ \bm{n_0} $                   & $ 1 $                                                                   & $ 1 $                                                                 & $ - $                                                               \\
                % $ \bm{n_{\text{reaccionan}}} $ & $ -2x $                                                                 & $ -x $                                                                & $ - $                                                               \\
                % $ \bm{n}_{\text{forman}} $     & $ - $                                                                   & $ - $                                                                 & $ 2x $                                                              \\
                % $ \bm{n}_{\text{eq}} $         & $ \underbrace{1-2x}_{{\color{red} 0.15}}  $                               & $ \underbrace{1-x}_{{\color{red} 0.575}}$                             & $ \underbrace{2x}_{{\color{red} \moles{0.85}}} $                    \\
                % \( \conc{\ \ }_{eq}\)          & $\underbrace{\dfrac{1-2x}{5}}_{{\color{blue}\concentracion[m]{0.03}}} $ & $\underbrace{\dfrac{1-x}{5}}_{{\color{blue}\concentracion[m]{0.11}}}$ & $\underbrace{\dfrac{x}{5}}_{{\color{blue}\concentracion[m]{0.17}}}$ \\
            % \end{tblr}
            % --- tabla nueva ---
            \begin{tabular}{@{}C{0.21\linewidth} @{\,}c@{\,} C{0.21\linewidth} @{\,}c@{\,} C{0.21\linewidth} @{\,}c@{\,} C{0.21\linewidth}@{}}
                \toprule
                \textbf{Ecuación} & : & \textbf{\ch{2 SO2\ (g)}} & \ch{+} & \textbf{\ch{O2 \ (g)}} & \ch{->} & \textbf{\ch{2 SO3\ (g)}} \\
                \midrule
                $ \bm{n_0} $ &  & $ 1 $ &  & $ 1 $ &  & $ - $ \\[3pt]
                $ \bm{n_{\text{reaccionan}}} $ &  & $ -2x $ &  & $ -x $ &  & $ - $ \\[3pt]
                $ \bm{n}_{\text{forman}} $ &  & $ - $ &  & $ - $ &  & $ 2x $ \\[3pt]
                $ \bm{n}_{\text{eq}} $ &  & $ \underbrace{1-2x}_{{\color{red} 0.15}}  $ &  & $ \underbrace{1-x}_{{\color{red} 0.575}}$ &  & $ \underbrace{2x}_{{\color{red} \moles{0.85}}} $ \\[3pt]
                \( \conc{\ \ }_{eq}\) &  & $\underbrace{\dfrac{1-2x}{5}}_{{\color{blue}\concentracion[m]{0.03}}} $ &  & $\underbrace{\dfrac{1-x}{5}}_{{\color{blue}\concentracion[m]{0.11}}}$ &  & $\underbrace{\dfrac{x}{5}}_{{\color{blue}\concentracion[m]{0.17}}}$ \\
                \bottomrule
            \end{tabular}
        \end{table}
        \begin{enumerate}[label=\letraminuscula]
\item Si sabemos que en el equilibrio tenemos \moles{0.15} de \ch{SO2} entonces:
\begin{equation*}
    1-2x = 0.15 \rightarrow x = \dfrac{1-0.15}{2} = \moles{0.425}
\end{equation*}
Ahora podemos calcular las concentraciones en equilibrio:
\begin{multicols}{2}
    \begin{itemize}
        \item \([\ch{SO2}] = \dfrac{0.15}{5} = \SI{0.03}{\molar}\).
        \item \([\ch{O2}] = \dfrac{0.85}{5} = \SI{0.115}{\molar}\).
        \item \([\ch{SO3}] = \dfrac{0.15}{5} = \SI{0.17}{\molar}\).
    \end{itemize}
\end{multicols}
\item Primero calculamos la \(k_c\) según:
\begin{equation*}
    k_c = \dfrac{\conc{SO3}^2}{\conc{SO2}^2\ \conc{O2}}
\end{equation*}
Sustituimos las concentraciones calculadas en el apartado anterior teniendo finalmente que:
\begin{equation*}
    k_c =\dfrac{0.17^2}{0.03^2\cdot 0.115} = 279.23\ \si{\per\molar}
\end{equation*}
Ahora calculamos la \(k_p\) según la relación:
\begin{equation*}
    k_p = k_c\ (R\ T)^{\Delta n} = 5.88\cdot (0.082\cdot 1273)^{-1} = \SI{2.67}{\per\atmospher}
\end{equation*}
\item Si aumentamos la presión del recipiente el equilibrio se desplazará hacia donde existan menos número de moles gaseosos, es decir aumentará la concentración de \ch{SO3} (equilibrio se desplaza hacia la derecha)
        \end{enumerate}
    \end{solu}

\end{exerciseEVAU}

\begin{exerciseEVAU}
    En un recipiente de \volumen[l]{10,0} se introduce una mezcla de \moles{4,0} de nitrógeno y \moles{12,0} moles de hidrógeno. Se eleva la temperatura hasta \temperatura[k]{1000} estableciéndose el equilibrio:
    \begin{equation*}
        \ch{N2\ (g) + 3 H2\ (g) <=> 2 NH3\ (g)}
    \end{equation*}
    En ese instante, se observa que hay \moles{0,8} de amoniaco en la mezcla gaseosa.
    Calcula:
    \begin{enumerate}[label=\letraminuscula]
\item La constante de concentraciones \(K_c\)
\item La constante de presiones \(K_p\)  y la presión total de la mezcla gaseosa en equilibrio.
\item Las presiones parciales de los componentes en el equilibrio.
    \end{enumerate}
    \datos{}: R = \ctegases

    \begin{solu}
%        \begin{figure}[H]
%            \centering
%            \includegraphics[scale=0.8]{images/imagesUD6EquilibrioQuimico/UD6_Problemas/UD6_EVAU_24A}
%            \includegraphics[scale=0.8]{images/imagesUD6EquilibrioQuimico/UD6_Problemas/UD6_EVAU_24B}
%        \end{figure}

        La reacción es \(\ch{N2 \gas{} + 3 H2 \gas{} <=> 2 NH3 \gas{}}\), con \(V = \SI{10.0}{\litre}\) y \(T = \SI{1000}{\kelvin}\). Como se han formado 0,8 mol de \ch{NH3}, y por estequiometría \(n(\ch{NH3}) = 2\xi\):

        \begin{equation*}
            \xi = \dfrac{0,8}{2} = 0,4 \ \text{mol}
        \end{equation*}

        \begin{equation*}
            n(\ch{N2})_{eq} = 4,0 - 0,4 = 3,6 \ \text{mol} \qquad n(\ch{H2})_{eq} = 12,0 - 3(0,4) = 10,8 \ \text{mol} \qquad n(\ch{NH3})_{eq} = 0,8 \ \text{mol}
        \end{equation*}

        \textbf{a) Constante \(K_c\)}

        \begin{equation*}
            [\ch{N2}] = \SI{0.36}{\Molar} \qquad [\ch{H2}] = \SI{1.08}{\Molar} \qquad [\ch{NH3}] = \SI{0.08}{\Molar}
        \end{equation*}

        \begin{equation*}
            K_c = \dfrac{[\ch{NH3}]^2}{[\ch{N2}][\ch{H2}]^3} = \dfrac{(0,08)^2}{0,36 \times (1,08)^3} = \num{0.0141}
        \end{equation*}

        \textbf{b) Constante \(K_p\) y presión total}

        Como \(\Delta n = 2-4=-2\):

        \begin{equation*}
            K_p = K_c\, (RT)^{-2} = \dfrac{0,0141}{(0,082 \times 1000)^2} = \num{2.10e-6}
        \end{equation*}

        \begin{equation*}
            n_T = 3,6+10,8+0,8 = 15,2 \ \text{mol} \qquad P_T = \dfrac{n_T\, R\, T}{V} = \dfrac{15,2 \times 0,082 \times 1000}{10,0} = \SI{124.6}{\atmospher}
        \end{equation*}

        \textbf{c) Presiones parciales}

        \begin{equation*}
            p_{\ch{N2}} = [\ch{N2}]\, RT = 0,36 \times 82 = \SI{29.5}{\atmospher} \qquad p_{\ch{H2}} = 1,08 \times 82 = \SI{88.6}{\atmospher} \qquad p_{\ch{NH3}} = 0,08 \times 82 = \SI{6.6}{\atmospher}
        \end{equation*}

    \end{solu}

\end{exerciseEVAU}

\begin{exerciseEVAU}
    El gas fosgeno, \ch{COCl2}, se obtiene con gran rendimiento mediante el proceso de equilibrio siguiente:
    \begin{equation*}
        \ch{CO\ (g) + Cl2\ (g) <=> COCl2\ (g)}
    \end{equation*}
    Si la concentración de \ch{CO} se incrementa mediante la adición de más \ch{CO} al sistema en equilibrio, sin cambio de temperatura, ¿cómo se verá afectada?:
    \begin{enumerate}[label=\letraminuscula]
\item ¿La concentración de \ch{COCl2} una vez que se haya alcanzado de nuevo el equilibrio?
\item ¿La constante de equilibrio de la reacción?
    \end{enumerate}

    \begin{solu}
        La constante de equilibrio referida a las concentraciones tendrá una expresión:
        \begin{equation*}
            k_c = \dfrac{\conc{COCl2}}{\conc{CO}\ \conc{Cl2}}
        \end{equation*}
        \begin{enumerate}[label=\letraminuscula]
\item Si la temperatura se mantiene constante, al aumentar la concentración de \ch{CO} el equilibrio se desplazará hacia los productos, es decir la concentración de fosgeno aumentará.
\item La constante de equilibrio no se ve afectada, ya que su valor solo se vería afectado si cambiamos la temperatura.
        \end{enumerate}
    \end{solu}

\end{exerciseEVAU}

\begin{exerciseEVAU}
    A \temperatura[c]{25}, la constante de equilibrio \(K_c\) de la reacción:
    \begin{equation*}
        \ch{N2O4\ (g)  <=> 2 NO2\ (g)}
    \end{equation*}
    vale \num{4,63e-3}.
    En un recipiente de \volumen[l]{2} se introducen \moles{1,6} de \ch{N2O4} a dicha temperatura.
    \begin{enumerate}[label=\letraminuscula]
\item ¿Cuáles serán las concentraciones de ambos gases en el equilibrio?
\item Calcula el valor de Kp y la presión total.
\item ¿Cómo afectaría a las concentraciones de los gases en el equilibrio una reducción del volumen del recipiente sin variar la temperatura? ¿Y la adición de un catalizador?
    \end{enumerate}

% \solution{a) \conc{NO2} = \concentracion[m]{0.06}; \conc{N2O4} = \concentracion[m]{0.77}; b) \(p_T = \SI{20.28}{\atmospher} \), \(k_p = 0.11\); c) La reducción del volumen desplaza el equilibrio hacia los reactivos y un catalizador solo aumenta la velocidad de la reacción}

    \begin{solu}
%        \begin{figure}[H]
%            \centering
%            \includegraphics[scale=0.8]{images/imagesUD6EquilibrioQuimico/UD6_Problemas/UD6_EVAU_26A}
%            \includegraphics[scale=0.8]{images/imagesUD6EquilibrioQuimico/UD6_Problemas/UD6_EVAU_26B}
%        \end{figure}

        La reacción es \(\ch{N2O4 \gas{} <=> 2 NO2 \gas{}}\), con \(V = \SI{2}{\litre}\) y \(T = \SI{298}{\kelvin}\). La concentración inicial es \([\ch{N2O4}]_0 = 1,6/2 = \SI{0.8}{\Molar}\).

        \textbf{a) Concentraciones en el equilibrio}

        \begin{equation*}
            K_c = \dfrac{(2x)^2}{0,8-x} = \num{4.63e-3} \ \Rightarrow \ 4x^2 + 0,00463\, x - 0,003704 = 0
        \end{equation*}

        Resolviendo la ecuación de segundo grado, \(x = 0,0299\ \si{\Molar}\), por lo que:

        \begin{equation*}
            [\ch{N2O4}]_{eq} = 0,8 - 0,0299 \approx \SI{0.770}{\Molar} \qquad [\ch{NO2}]_{eq} = 2(0,0299) \approx \SI{0.060}{\Molar}
        \end{equation*}

        \textbf{b) \(K_p\) y presión total}

        Como \(\Delta n = 2-1=1\):

        \begin{equation*}
            K_p = K_c\, (RT) = 4,63 \times 10^{-3} \times (0,082 \times 298) = \num{0.113}
        \end{equation*}

        \begin{equation*}
            n_T = ([\ch{N2O4}]+[\ch{NO2}]) \times V = (0,770+0,060) \times 2 = 1,66 \ \text{mol}
        \end{equation*}

        \begin{equation*}
            P_T = \dfrac{n_T\, R\, T}{V} = \dfrac{1,66 \times 0,082 \times 298}{2} = \SI{20.28}{\atmospher}
        \end{equation*}

        \textbf{c) Efecto de reducir el volumen y de un catalizador}

        Al \textbf{reducir el volumen} (aumentar la presión) manteniendo la temperatura, el equilibrio se desplaza hacia el lado con menos moles gaseosos, es decir, hacia el \ch{N2O4} (reactivos), disminuyendo la concentración de \ch{NO2}.

        La \textbf{adición de un catalizador} no modifica las concentraciones de equilibrio ni el valor de \(K_c\), ya que un catalizador solo aumenta la velocidad con la que se alcanza el equilibrio, sin afectar a su posición.

    \end{solu}

\end{exerciseEVAU}

\begin{exerciseEVAU}
    \moles{2.3e-2} moles de gas fosgeno (\ch{COCl2}) puro se introdujeron en un reactor de \volumen[l]{1,50}, calentándose entonces éste hasta alcanzar los \temperatura[k]{800} de temperatura. Alcanzado el equilibrio, la presión parcial de \ch{CO} fue \presion[atm]{0,497}. Para el equilibrio
    \begin{equation*}
        \ch{COCl2\ (g) <=> CO\ (g) + Cl2\ (g)}
    \end{equation*}
    calcula:
    \begin{enumerate}[label=\letraminuscula]
\item La constante de presiones \(K_p\).
\item El número total de moles en el equilibrio.
\item El grado de disociación del gas fosgeno
    \end{enumerate}
    \datos{}: R = \ctegases

% \solution{a) \(k_p = 0.487\), b) \(n_T = \moles{0.0344}\), c) \(\alpha = 0.5\)}

    \begin{solu}
%        \begin{figure}[H]
%            \centering
%            \includegraphics[scale=0.8]{images/imagesUD6EquilibrioQuimico/UD6_Problemas/UD6_EVAU_27}
%        \end{figure}

        La reacción es \(\ch{COCl2 \gas{} <=> CO \gas{} + Cl2 \gas{}}\), con \(n_0 = \num{2.3e-2}\) mol, \(V = \SI{1.50}{\litre}\) y \(T = \SI{800}{\kelvin}\). Como \([\ch{CO}] = [\ch{Cl2}]\) (parten de cero en igual proporción), y \(p_{\ch{CO}} = \SI{0.497}{\atmospher}\):

        \begin{equation*}
            n(\ch{CO})_{eq} = n(\ch{Cl2})_{eq} = \dfrac{p_{\ch{CO}}\, V}{R\, T} = \dfrac{0,497 \times 1,50}{0,082 \times 800} = \num{1.136e-2} \ \text{mol}
        \end{equation*}

        \begin{equation*}
            n(\ch{COCl2})_{eq} = 0,023 - 0,01136 = \num{1.164e-2} \ \text{mol}
        \end{equation*}

        \begin{equation*}
            p_{\ch{COCl2}} = \dfrac{n(\ch{COCl2})\, R\, T}{V} = \dfrac{0,01164 \times 0,082 \times 800}{1,50} = \SI{0.509}{\atmospher}
        \end{equation*}

        \textbf{a) Constante \(K_p\)}

        \begin{equation*}
            K_p = \dfrac{p_{\ch{CO}}\, p_{\ch{Cl2}}}{p_{\ch{COCl2}}} = \dfrac{(0,497)^2}{0,509} = \num{0.485}
        \end{equation*}

        \textbf{b) Moles totales en el equilibrio}

        \begin{equation*}
            n_T = 0,01164 + 0,01136 + 0,01136 = \num{0.0344} \ \text{mol}
        \end{equation*}

        \textbf{c) Grado de disociación}

        \begin{equation*}
            \alpha = \dfrac{n(\ch{CO})_{eq}}{n_0} = \dfrac{0,01136}{0,023} = 0,494 \approx 0,5 \ (50\ \%)
        \end{equation*}

    \end{solu}


\end{exerciseEVAU}

\begin{exerciseEVAU}
    \moles{2,50} de \ch{NOCl} puro se introdujeron en un reactor de \volumen[l]{1,50} a \temperatura[c]{400}. Una vez alcanzado el equilibrio, se comprobó que el 28.0\% del \ch{NOCl} inicial se había disociado según la reacción:
    \begin{equation*}
        \ch{2 NOCl\ (g) <=> 2 NO\ (g) + Cl2\ (g)}
    \end{equation*}
    Calcula:
    \begin{enumerate}[label=\letraminuscula]
\item La constante de concentraciones \(K_c\)  para este equilibrio.
\item La presión total en el equilibrio.
\item La constante de presiones \(K_p\)  para este equilibrio.
    \end{enumerate}
    \datos{}: R = \ctegases

% \solution{a) \(k_c=0.0357\); b) \(p_T = \SI{104.83}{\atmospher} \); c) \(k_p = 1.95\) }

    \begin{solu}
%        \begin{figure}[H]
%            \centering
%            \includegraphics[scale=0.8]{images/imagesUD6EquilibrioQuimico/UD6_Problemas/UD6_EVAU_28}
%        \end{figure}

        La reacción es \(\ch{2 NOCl \gas{} <=> 2 NO \gas{} + Cl2 \gas{}}\), con \(n_0(\ch{NOCl}) = 2,50\) mol, \(V = \SI{1.50}{\litre}\) y \(T = \SI{673}{\kelvin}\) (\(400+273\)).

        Como se disocia el 28,0\,\% del \ch{NOCl} inicial:

        \begin{equation*}
            n(\ch{NOCl})_{disociado} = 0,28 \times 2,50 = 0,70 \ \text{mol}
        \end{equation*}

        \begin{equation*}
            n(\ch{NOCl})_{eq} = 2,50 - 0,70 = 1,80 \ \text{mol} \qquad n(\ch{NO})_{eq} = 0,70 \ \text{mol} \qquad n(\ch{Cl2})_{eq} = \dfrac{0,70}{2} = 0,35 \ \text{mol}
        \end{equation*}

        \textbf{a) Constante \(K_c\)}

        \begin{equation*}
            [\ch{NOCl}] = \dfrac{1,80}{1,50} = \SI{1.20}{\Molar} \qquad [\ch{NO}] = \dfrac{0,70}{1,50} = \SI{0.467}{\Molar} \qquad [\ch{Cl2}] = \dfrac{0,35}{1,50} = \SI{0.233}{\Molar}
        \end{equation*}

        \begin{equation*}
            K_c = \dfrac{[\ch{NO}]^2[\ch{Cl2}]}{[\ch{NOCl}]^2} = \dfrac{(0,467)^2 \times 0,233}{(1,20)^2} = \num{0.0353}
        \end{equation*}

        \textbf{b) Presión total}

        \begin{equation*}
            n_T = 1,80+0,70+0,35 = 2,85 \ \text{mol}
        \end{equation*}

        \begin{equation*}
            P_T = \dfrac{n_T\, R\, T}{V} = \dfrac{2,85 \times 0,082 \times 673}{1,50} = \SI{104.86}{\atmospher}
        \end{equation*}

        \textbf{c) Constante \(K_p\)}

        Como \(\Delta n = 3-2=1\):

        \begin{equation*}
            K_p = K_c\, (RT) = 0,0353 \times (0,082 \times 673) = \num{1.95}
        \end{equation*}

    \end{solu}

\end{exerciseEVAU}

\begin{exerciseEVAU}
    En el equilibrio:
    \begin{equation*}
        \ch{PCl5\ (g) <=> PCl3\ (g) + Cl2\ (g)},\; \Delta H^0 = \energia[kjmol]{92,5}
    \end{equation*}
    predice razonadamente el cambio que se producirá cuando:
    \begin{enumerate}[label=\letraminuscula]
\item Se eleva la temperatura.
\item Se añade más gas cloro a la mezcla de reacción.
\item Se extrae \ch{PCl3} de la mezcla de reacción.
\item Se incrementa la presión total del sistema
    \end{enumerate}

    \begin{solu}
%        \begin{figure}[H]
%            \centering
%            \includegraphics[scale=0.8]{images/imagesUD6EquilibrioQuimico/UD6_Problemas/UD6_EVAU_29}
%        \end{figure}

        La reacción \(\ch{PCl5 \gas{} <=> PCl3 \gas{} + Cl2 \gas{}}\) tiene \(\Delta H^0 = \energia[kjmol]{92.5} > 0\), es decir, es \textbf{endotérmica} en el sentido de la disociación (hacia la derecha).

        \textbf{a) Aumento de la temperatura}

        Un aumento de temperatura desplaza el equilibrio en el sentido que tiende a absorber el exceso de energía, es decir, en el sentido endotérmico. El equilibrio se desplaza \textbf{hacia la derecha}, aumentando la disociación del \ch{PCl5}.

        \textbf{b) Adición de más \ch{Cl2}}

        Al añadir \ch{Cl2}, aumenta la concentración de uno de los productos, por lo que, según el principio de Le Châtelier, el equilibrio se desplaza \textbf{hacia la izquierda}, consumiendo parte del \ch{Cl2} añadido y formando más \ch{PCl5}.

        \textbf{c) Extracción de \ch{PCl3}}

        Al extraer \ch{PCl3} de la mezcla, disminuye su concentración, por lo que el equilibrio se desplaza \textbf{hacia la derecha} para reponer parte del \ch{PCl3} extraído, aumentando la disociación del \ch{PCl5}.

        \textbf{d) Aumento de la presión total}

        Un aumento de la presión total del sistema (por ejemplo, al reducir el volumen) desplaza el equilibrio hacia el lado con menor número de moles gaseosos. Como hay 1 mol de gas en reactivos frente a 2 mol en productos, el equilibrio se desplaza \textbf{hacia la izquierda}, disminuyendo la disociación del \ch{PCl5}.

    \end{solu}
\end{exerciseEVAU}

